%!TEX root = forallxyyc.tex
\part{Lógica de primeira ordem}
\label{ch.FOL}
\addtocontents{toc}{\protect\mbox{}\protect\hrulefill\par}
\chapter{Blocos fundamentais da LPO}\label{s:FOLBuildingBlocks}

\section{A necessidade de decompor sentenças}
Considere o argumento a seguir, que é obviamente válido em inglês:
\begin{earg}
\label{willard1}
\item[] Willard é um lógico.
\item[] Todos os lógicos usam chapéus engraçados.
\item[\texttherefore] Willard usa um chapéu engraçado.
\end{earg}
Para simbolizá-lo na LVF, poderíamos fornecer uma chave de simbolização:
\begin{ekey}
\item[L] Willard é um lógico.
\item[A] Todos os lógicos usam chapéus engraçados.
\item[F] Willard usa um chapéu engraçado.
\end{ekey}
E o próprio argumento se torna:
$$L, A \therefore F$$
Mas o teste por tabela-verdade agora indicará que ele é \emph{inválido}. O que deu errado?

O problema não é termos cometido um erro ao simbolizar o argumento. Essa é a melhor simbolização que podemos fornecer \emph{na LVF}. O problema está na própria LVF. \textquotedblleft Todos os lógicos usam chapéus engraçados\textquotedblright{} trata tanto dos lógicos quanto do uso de chapéus. Ao não conservar essa estrutura em nossa simbolização, perdemos a conexão entre o fato de Willard ser um lógico e o fato de usar um chapéu.

As unidades básicas da LVF são letras sentenciais, e a LVF não consegue decompor essas letras. Para simbolizar argumentos como o anterior, teremos de desenvolver uma nova linguagem lógica que nos permita \emph{decompor o átomo}. Chamaremos essa linguagem de linguagem da \emph{lógica de primeira ordem}, ou \emph{LPO}.\footnote{A lógica de primeira ordem também é frequentemente chamada de lógica (clássica) de predicados ou lógica quantificacional.}

Os detalhes da LPO serão explicados ao longo deste capítulo, mas eis a ideia básica para decompor o átomo.

Primeiro, temos \emph{nomes}. Na LPO, indicamos os nomes com letras minúsculas em itálico. Por exemplo, poderíamos fazer \textquotedblleft $b$\textquotedblright{} representar Bertie ou fazer \textquotedblleft $i$\textquotedblright{} representar Willard.

Segundo, temos predicados. Predicados em inglês são expressões como \textquotedblleft\blank\ é um cão\textquotedblright{} ou \textquotedblleft\blank\ é um lógico\textquotedblright{}. Essas expressões não são sentenças completas por si mesmas. Para construir uma sentença completa, precisamos preencher a lacuna. Precisamos dizer algo como \textquotedblleft Bertie é um cão\textquotedblright{} ou \textquotedblleft Willard é um lógico\textquotedblright{}. Na LPO, indicamos predicados com letras maiúsculas em itálico. Por exemplo, poderíamos fazer o predicado da LPO \textquotedblleft $D$\textquotedblright{} simbolizar o predicado em inglês \textquotedblleft\blank\ é um cão\textquotedblright{}. Então, a expressão \textquotedblleft $\atom{D}{b}$\textquotedblright{} será uma sentença da LPO que simboliza a sentença em inglês \textquotedblleft Bertie é um cão\textquotedblright{}. Da mesma forma, poderíamos fazer o predicado da LPO \textquotedblleft $L$\textquotedblright{} simbolizar o predicado em inglês \textquotedblleft\blank\ é um lógico\textquotedblright{}. Então, a expressão \textquotedblleft $\atom{L}{i}$\textquotedblright{} simbolizará a sentença em inglês \textquotedblleft Willard é um lógico\textquotedblright{}.

Terceiro, temos quantificadores. Por exemplo, \textquotedblleft $\exists$\textquotedblright{} transmitirá aproximadamente a ideia de \textquotedblleft existe pelo menos um\ldots\textquotedblright{}. Assim, poderíamos simbolizar a sentença em inglês \textquotedblleft existe um cão\textquotedblright{} pela sentença da LPO \textquotedblleft $\exists x\, \atom{D}{x}$\textquotedblright{}, que poderíamos ler em voz alta como \textquotedblleft existe pelo menos uma coisa, $x$, tal que $x$ é um cão\textquotedblright{}.

Essa é a ideia geral, mas a LPO é significativamente mais sutil que a LVF; avançaremos devagar.

\section{Nomes}
Em inglês, um \emph{termo singular} é uma palavra ou expressão que pretende referir-se a uma pessoa, lugar ou coisa \emph{específica}. A palavra \textquotedblleft cão\textquotedblright{} não é um termo singular, porque há muitos cães. O nome \textquotedblleft Bertie\textquotedblright{} é um termo singular porque se refere a um terrier específico. Da mesma forma, a expressão \textquotedblleft cão de Philip\textquotedblright{} é um termo singular porque também se refere a um pequeno terrier específico.

\emph{Nomes próprios} são um tipo particularmente importante de termo singular. São expressões que identificam indivíduos sem descrevê-los. O nome \textquotedblleft Emerson\textquotedblright{} é um nome próprio, e o nome, por si só, não informa nada sobre Emerson. É claro que alguns nomes são tradicionalmente dados a meninos e outros, a meninas. Se \textquotedblleft Hilary\textquotedblright{} for usado como termo singular, você talvez suponha que se refere a uma mulher. No entanto, essa suposição pode estar errada. De fato, o nome não significa necessariamente que a pessoa referida seja sequer uma pessoa: Hilary pode ser uma girafa, tanto quanto você poderia saber apenas pelo nome.

Na LPO, nossos \define{nomes} são letras minúsculas de \textquotedblleft $a$\textquotedblright{} a \textquotedblleft $r$\textquotedblright{}. Podemos acrescentar subscritos se quisermos usar alguma letra mais de uma vez. Eis alguns termos singulares na LPO:
	$$a,b,c,\ldots, r, a_1, f_{32}, j_{390}, m_{12}$$
Eles devem ser entendidos como nomes próprios em inglês, mas com uma diferença. \textquotedblleft Tim Button\textquotedblright{} é um nome próprio, mas há várias pessoas com esse nome. Vivemos com esse tipo de ambiguidade em inglês, permitindo que o contexto individualize o fato de que \textquotedblleft Tim Button\textquotedblright{} se refere ao autor deste livro, e não a outro Tim. Na LPO, não toleramos nenhuma ambiguidade desse tipo. Cada nome deve designar \emph{exatamente} uma coisa. Entretanto, assim como em inglês, nomes diferentes podem designar a mesma coisa, e pode haver coisas que não sejam designadas por nenhum nome.

\newglossaryentry{name}{
  name = nome,
  description = {Um símbolo da LPO usado para designar um objeto do \gls{domain}}
  }

Assim como na LVF, podemos fornecer chaves de simbolização. Elas indicam, temporariamente, o que um nome designará. Assim, poderíamos oferecer:
	\begin{ekey}
		\item[e] Elsa
		\item[g] Gregor
		\item[m] Marybeth
	\end{ekey}


\section{Predicados}
Os predicados mais simples expressam propriedades de indivíduos. São coisas que podemos afirmar sobre um objeto. Eis alguns exemplos de predicados em inglês:
	\begin{itemize}
		\item \blank\ é um cão
		\item \blank\ é membro do Monty Python
		\item Um piano caiu sobre \blank
	\end{itemize}
Em geral, podemos pensar nos predicados como coisas que se combinam com termos singulares para formar sentenças. Ao contrário, podemos começar com sentenças e formar predicados a partir delas removendo termos. Considere a sentença \textquotedblleft Vinnie pegou emprestado o carro da família de Nunzio\textquotedblright{}. Ao remover um termo singular, podemos obter qualquer um destes três predicados diferentes:
	\begin{itemize}
		\item \blank\ pegou emprestado o carro da família de Nunzio\\
		\item Vinnie pegou emprestado \blank\ de Nunzio\\
		\item Vinnie pegou emprestado o carro da família de \blank
	\end{itemize}
Na LPO, \define{predicados} são letras maiúsculas de \textquotedblleft $A$\textquotedblright{} a \textquotedblleft $Z$\textquotedblright{}, com ou sem subscritos. Poderíamos escrever uma chave de simbolização para predicados da seguinte forma:
	\begin{ekey}
		\item[\atom{A}{x}] \gap{x} está zangado
		\item[\atom{H}{x}] \gap{x} está feliz
%		\item[\atom{T_1}{x,y}] \gap{x} é tão alto quanto ou mais alto que \gap{y}
%		\item[\atom{T_2}{x,y}] \gap{x} é tão resistente quanto ou mais resistente que \gap{y}
%		\item[\atom{B}{x,y,z}] \gap{y} está entre \gap{x} e \gap{z}
	\end{ekey}
        (Por que os subscritos nas lacunas? Voltaremos a isso em \cref{s:MultipleGenerality}.)

\newglossaryentry{predicate}{
  name = predicado,
  description = {Um símbolo da LPO usado para simbolizar uma propriedade ou relação}
}

Se combinarmos nossas duas chaves de simbolização, poderemos começar a simbolizar algumas sentenças em inglês que usam esses nomes e predicados em conjunto. Por exemplo, considere as sentenças em inglês:
	\begin{enumerate}
		\item\label{terms1} Elsa está zangada.
		\item\label{terms2a} Gregor e Marybeth estão zangados.
		\item\label{terms2} Se Elsa está zangada, Gregor e Marybeth também estão.
	\end{enumerate}
\Cref*{terms1} é direta: nós a simbolizamos por \textquotedblleft $\atom{A}{e}$\textquotedblright{}.

\Cref*{terms2a} é uma conjunção de duas sentenças mais simples. As sentenças simples podem ser simbolizadas apenas por \textquotedblleft $\atom{A}{g}$\textquotedblright{} e \textquotedblleft $\atom{A}{m}$\textquotedblright{}. Então, recorremos aos nossos recursos da LVF e simbolizamos a sentença inteira por \textquotedblleft $\atom{A}{g} \eand \atom{A}{m}$\textquotedblright{}. Isso ilustra um ponto importante: a LPO possui todos os conectivos verofuncionais da LVF.

\Cref*{terms2} é um condicional cujo antecedente é \cref*{terms1} e cujo consequente é \cref*{terms2a}; portanto, podemos simbolizá-la por \textquotedblleft $\atom{A}{e} \eif (\atom{A}{g} \eand \atom{A}{m})$\textquotedblright{}.

\section{Quantificadores}
Estamos prontos para introduzir os quantificadores. Considere estas sentenças:
	\begin{enumerate}
		\item\label{q.a} Todos estão felizes.
%		\item\label{q.ac} Todos são pelo menos tão resistentes quanto Elsa.
		\item\label{q.e} Alguém está zangado.
	\end{enumerate}
Pode ser tentador simbolizar \cref*{q.a} como \textquotedblleft $\atom{H}{e} \eand \atom{H}{g} \eand \atom{H}{m}$\textquotedblright{}. No entanto, isso diria apenas que Elsa, Gregor e Marybeth estão felizes. Queremos dizer que \emph{todos} estão felizes, inclusive aqueles que não têm nome. Para fazer isso, introduzimos o símbolo \textquotedblleft $\forall$\textquotedblright{}. Ele é chamado de \define{quantificador universal}.

\newglossaryentry{universal quantifier}{
  name = quantificador universal,
  description = {O símbolo $\forall$ da LPO usado para simbolizar generalidade; $\forall x\, \atom{F}{x}$ é verdadeiro \ifeff{} todo membro do domínio é~$F$}
}

Um quantificador deve sempre ser seguido por uma \define{vari\'avel}. Na LPO, variáveis são letras minúsculas em itálico de \textquotedblleft $s$\textquotedblright{} a \textquotedblleft $z$\textquotedblright{}, com ou sem subscritos. Assim, poderíamos simbolizar \cref*{q.a} como \textquotedblleft $\forall x\, \atom{H}{x}$\textquotedblright{}. A variável \textquotedblleft $x$\textquotedblright{} funciona como uma espécie de espaço reservado. A expressão \textquotedblleft $\forall x$\textquotedblright{} significa intuitivamente que você pode escolher qualquer pessoa e colocá-la no lugar de \textquotedblleft $x$\textquotedblright{}. A expressão seguinte, \textquotedblleft $\atom{H}{x}$\textquotedblright{}, indica que essa coisa que você escolheu é feliz.

\newglossaryentry{variable}{
  name = variável,
  description = {Um símbolo da LPO usado depois dos quantificadores e como espaço reservado em fórmulas atômicas; letras minúsculas entre $s$ e~$z$}
}


Deve-se observar que não há nenhuma razão especial para usar \textquotedblleft $x$\textquotedblright{} em vez de outra variável. As sentenças \textquotedblleft $\forall x\, \atom{H}{x}$\textquotedblright{}, \textquotedblleft $\forall y\, \atom{H}{y}$\textquotedblright{}, \textquotedblleft $\forall z\, \atom{H}{z}$\textquotedblright{} e \textquotedblleft $\forall x_5\, \atom{H}{x_5}$\textquotedblright{} usam variáveis diferentes, mas todas são logicamente equivalentes.

Para simbolizar \cref{q.e}, introduzimos outro símbolo novo: o \define{quantificador existencial}, \textquotedblleft $\exists$\textquotedblright{}. Assim como o quantificador universal, o quantificador existencial exige uma variável. \Cref*{q.e} pode ser simbolizada por \textquotedblleft $\exists x\,\atom{A}{x}$\textquotedblright{}. Enquanto \textquotedblleft $\forall x\,\atom{A}{x}$\textquotedblright{} é naturalmente lida como \textquotedblleft para todo $x$, $x$ está zangado\textquotedblright{}, \textquotedblleft $\exists x\,\atom{A}{x}$\textquotedblright{} é naturalmente lida como \textquotedblleft existe alguma coisa, $x$, tal que $x$ está zangado\textquotedblright{}. Mais uma vez, a variável é uma espécie de espaço reservado; poderíamos ter simbolizado \cref*{q.e} com a mesma facilidade por \textquotedblleft $\exists z\, \atom{A}{z}$\textquotedblright{}, \textquotedblleft $\exists w_{256}\, \atom{A}{w_{256}}$\textquotedblright{}, ou qualquer outra variável.

\newglossaryentry{existential quantifier}{
  name = quantificador existencial,
  description = {O símbolo $\exists$ da LPO usado para simbolizar a existência; $\exists x\, \atom{F}{x}$ é verdadeiro \ifeff{} pelo menos um membro do domínio é~$F$}
}

Alguns exemplos ajudarão. Considere estas outras sentenças:
	\begin{enumerate}
		\item\label{q.ne} Ninguém está zangado.
		\item\label{q.en} Há alguém que não está feliz.
		\item\label{q.na} Nem todos estão felizes.
	\end{enumerate}
\Cref*{q.ne} pode ser parafraseada como \textquotedblleft não é o caso de que alguém está zangado\textquotedblright{}. Podemos então simbolizá-la usando a negação e um quantificador existencial: \textquotedblleft $\enot \exists x\,\atom{A}{x}$\textquotedblright{}. No entanto, \cref*{q.ne} também poderia ser parafraseada como \textquotedblleft todas as pessoas não estão zangadas\textquotedblright{}. Com isso em mente, ela pode ser simbolizada usando a negação e um quantificador universal: \textquotedblleft $\forall x\, \enot \atom{A}{x}$\textquotedblright{}. Ambas as simbolizações são aceitáveis. De fato, veremos que, em geral, $\forall x\, \enot\metav{A}$ é logicamente equivalente a $\enot\exists x\,\metav{A}$. (Observe que aqui retornamos à prática de usar $\metav{A}$ como metavariável, de \cref{s:UseMention}.) Simbolizar uma sentença de uma maneira ou de outra pode parecer mais \textquotedblleft natural\textquotedblright{} em alguns contextos, mas não passa de uma questão de gosto.

\Cref*{q.en} é mais naturalmente parafraseada como \textquotedblleft existe algum $x$ tal que $x$ não está feliz\textquotedblright{}. Isso se torna \textquotedblleft $\exists x\, \enot \atom{H}{x}$\textquotedblright{}. É claro que poderíamos igualmente ter escrito \textquotedblleft $\enot\forall x\,\atom{H}{x}$\textquotedblright{}, que leríamos naturalmente como \textquotedblleft não é o caso de que todos estão felizes\textquotedblright{}. Essa também seria uma simbolização perfeitamente adequada de \cref*{q.na}.


\section{Domínios}
Dada a chave de simbolização que temos usado, \textquotedblleft $\forall x\,\atom{H}{x}$\textquotedblright{} simboliza \textquotedblleft todos estão felizes\textquotedblright{}. Quem está incluído nesse \emph{todos}? Quando usamos sentenças como essa em inglês, normalmente não queremos dizer todas as pessoas atualmente vivas na Terra. Certamente não queremos dizer todas as pessoas que já viveram ou viverão. Normalmente queremos dizer algo mais modesto: todas as pessoas que estão no prédio, todas as pessoas matriculadas na aula de balé ou algo semelhante.

Para eliminar essa ambiguidade, precisaremos especificar um \define{dom\'inio}. O domínio é a coleção de coisas sobre as quais estamos falando. Assim, se quisermos falar sobre pessoas em Chicago, definimos o domínio como pessoas em Chicago. Escrevemos isso no início da chave de simbolização, assim:
	\begin{ekey}
		\item[\text{domínio}] pessoas em Chicago
	\end{ekey}
Os quantificadores \emph{varrem} o domínio. Dado esse domínio, \textquotedblleft $\forall x$\textquotedblright{} deve ser lido aproximadamente como \textquotedblleft toda pessoa em Chicago é tal que\ldots\textquotedblright{} e \textquotedblleft $\exists x$\textquotedblright{} deve ser lido aproximadamente como \textquotedblleft alguma pessoa em Chicago é tal que\ldots\textquotedblright{}.

\newglossaryentry{domain}{
  name = domínio,
  description = {A coleção de objetos assumida para uma simbolização na LPO ou que determina o alcance dos quantificadores em uma \gls{interpretation}}
}

Na LPO, o domínio deve sempre incluir pelo menos uma coisa. Além disso, em inglês podemos inferir legitimamente \textquotedblleft alguma coisa está zangada\textquotedblright{} de \textquotedblleft Gregor está zangado\textquotedblright{}. Portanto, na LPO, queremos poder inferir \textquotedblleft $\exists x\,\atom{A}{x}$\textquotedblright{} de \textquotedblleft $\atom{A}{g}$\textquotedblright{}. Assim, exigiremos que cada nome designe exatamente uma coisa no domínio. Se quisermos nomear pessoas em lugares além de Chicago, precisaremos incluí-las no domínio.
	\factoidbox{
		Um domínio deve ter \emph{pelo menos} um membro. Cada nome deve designar \emph{exatamente} um membro do domínio, mas um membro do domínio pode ser designado por um nome, por muitos nomes ou por nenhum nome.
	}

Mesmo permitir um domínio com apenas um membro pode produzir alguns resultados estranhos. Suponha que tenhamos esta chave de simbolização:
\begin{ekey}
\item[\text{domínio}] a Torre Eiffel
\item[\atom{P}{x}] \gap{x} está em Paris.
\end{ekey}
A sentença $\forall x\,\atom{P}{x}$ poderia ser parafraseada em inglês como \textquotedblleft tudo está em Paris\textquotedblright{}. Contudo, isso seria enganoso. Ela significa que tudo \emph{no domínio} está em Paris. Esse domínio contém apenas a Torre Eiffel; portanto, com essa chave de simbolização, $\forall x\,\atom{P}{x}$ significa simplesmente que a Torre Eiffel está em Paris.

\section{Termos não referentes}

Na LPO, cada nome deve designar exatamente um membro do domínio. Um nome não pode referir-se a mais de uma coisa — é um termo \emph{singular}. Cada nome ainda deve designar \emph{alguma coisa}. Isso está relacionado a um problema filosófico clássico: o chamado problema dos termos não referentes.

Os filósofos medievais normalmente usavam sentenças sobre a \emph{quimera} para exemplificar esse problema. A quimera é uma criatura mitológica; ela não existe de fato. Considere estas duas sentenças:
\begin{enumerate}
\item\label{chimera1} A quimera está zangada.
\item\label{chimera2} A quimera não está zangada.
\end{enumerate}
É tentador simplesmente definir um nome como significando \textquotedblleft quimera\textquotedblright{}. A chave de simbolização teria esta aparência:
\begin{ekey}
\item[\text{domínio}] criaturas na Terra
\item[\atom{A}{x}] \gap{x} está zangado.
\item[c] quimera
\end{ekey}
Poderíamos então simbolizar \cref*{chimera1} como $\atom{A}{c}$ e \cref*{chimera2} como $\enot \atom{A}{c}$.

Surgirão problemas quando perguntarmos se essas sentenças são verdadeiras ou falsas.

Uma opção é dizer que \cref*{chimera1} não é verdadeira, porque não existe nenhuma quimera. Se \cref*{chimera1} for falsa porque fala de algo inexistente, então \cref*{chimera2} será falsa pelo mesmo motivo. No entanto, isso significaria que tanto $\atom{A}{c}$ quanto $\enot \atom{A}{c}$ seriam falsas. Dadas as condições de verdade da negação, isso não pode acontecer.

Como não podemos dizer que ambas são falsas, o que devemos fazer? Outra opção é dizer que \cref*{chimera1} é \emph{sem sentido} porque fala de algo inexistente. Assim, $\atom{A}{c}$ seria uma expressão com sentido na LPO para algumas interpretações, mas não para outras. No entanto, isso tornaria nossa linguagem formal refém de interpretações particulares. Como estamos interessados na forma lógica, queremos considerar a força lógica de uma sentença como $\atom{A}{c}$ independentemente de qualquer interpretação particular. Se $\atom{A}{c}$ tivesse sentido em alguns casos e não tivesse sentido em outros, não poderíamos fazer isso.

Esse é o \emph{problema dos termos não referentes}, e voltaremos a ele mais adiante (ver \cref{subsec.defdesc}). O ponto importante por enquanto é que cada nome da LPO \emph{deve} referir-se a alguma coisa no domínio, embora o domínio possa conter quaisquer coisas que desejemos. Se quisermos simbolizar argumentos sobre criaturas mitológicas, teremos de definir um domínio que as inclua. Essa opção é importante se quisermos considerar a lógica de histórias. Podemos simbolizar uma sentença como \textquotedblleft Sherlock Holmes morou em 221B Baker Street\textquotedblright{} incluindo personagens fictícios como Sherlock Holmes em nosso domínio.

\chapter{Sentenças com um quantificador}
\label{s:MoreMonadic}

Agora temos todas as peças da LPO. Simbolizar sentenças mais complicadas é apenas uma questão de saber como combinar predicados, nomes, quantificadores e conectivos. Há uma certa habilidade nisso, e não há substituto para a prática.

%\section{Dealing with syncategorematic adjectives}
%When we encounter a sentence like
%	\begin{enumerate}
%		\item\label{syn1} Herbie is a white car
%	\end{enumerate}
%We can paraphrase this as `Herbie is white and Herbie is a car'. We can then use a symbolization key like:
%	\begin{ekey}
%		\item[\atom{W}{x}] \gap{x} is white
%		\item[\atom{C}{x}] \gap{x} is a car
%		\item[h] Herbie
%	\end{ekey}
%This allows us to symbolize \cref*{syn1} as `$\atom{W}{h} \eand \atom{C}{h}$'. But now consider:
%	\begin{enumerate}
%		\item\label{syn2} Damon Stoudamire is a short basketball player.
%		\item\label{syn3} Damon Stoudamire is a man.
%		\item\label{syn4} Damon Stoudamire is a short man.
%	\end{enumerate}
%Following the case of Herbie, we might try to use a symbolization key like:
%	\begin{ekey}
%		\item[\atom{S}{x}] \gap{x} is short
%		\item[\atom{B}{x}] \gap{x} is a basketball player
%		\item[\atom{M}{x}] \gap{x} is a man
%		\item[d] Damon Stoudamire
%	\end{ekey}
%Then we would symbolize \cref*{syn2} with `$\atom{S}{d} \eand \atom{B}{d}$', \cref*{syn3} with `$\atom{M}{d}$' and \cref*{syn4} with `$\atom{S}{d} \eand \atom{M}{d}$', but that would be a terrible mistake! This now suggests that \cref{syn2,syn3} together \emph{entail} \cref*{syn4}, but they do not. Standing at  5'10'', Damon Stoudamire is one of the shortest professional basketball players of all time, but he is nevertheless an averagely-tall man. The point is that \cref*{syn2} says that Damon is short \emph{qua} basketball player, even though he is of average height \emph{qua} man. So you will need to symbolize `\blank\ is a short basketball player' and `\blank\ is a short man' using completely different predicates.

%Similar examples abound. All politicians are people, but some good politicians are (arguably) bad people. Someone might be an incompetent statesman, but a competent individual. And so it goes. The moral is: when you see two adjectives in a row, you need to ask yourself carefully whether they can be treated as a conjunction or not.


\section{Expressões quantificadoras comuns}
Considere estas sentenças:
	\begin{enumerate}
\item\label{quan1} Toda moeda no meu bolso é uma moeda de vinte e cinco centavos.
\item\label{quan2} Alguma moeda sobre a mesa é uma moeda de dez centavos.
\item\label{quan3} Nem todas as moedas sobre a mesa são moedas de dez centavos.
\item\label{quan4} Nenhuma das moedas no meu bolso é uma moeda de dez centavos.
	\end{enumerate}
Ao fornecer uma chave de simbolização, precisamos especificar um domínio. Como estamos falando das moedas no meu bolso e sobre a mesa, o domínio deve conter pelo menos todas essas moedas. Como não estamos falando de nada além de moedas, fazemos com que o domínio seja formado por todas as moedas. Como não estamos falando de moedas específicas, não precisamos lidar com nomes. Eis, então, nossa chave:
	\begin{ekey}
\item[\text{domínio}] todas as moedas
\item[\atom{P}{x}] \gap{x} está no meu bolso
\item[\atom{T}{x}] \gap{x} está sobre a mesa
\item[\atom{Q}{x}] \gap{x} é uma moeda de vinte e cinco centavos
\item[\atom{D}{x}] \gap{x} é uma moeda de dez centavos
	\end{ekey}
\Cref*{quan1} é naturalmente simbolizada usando um quantificador universal. O quantificador universal diz algo sobre tudo no domínio, não apenas sobre as moedas no meu bolso. \Cref{quan1} pode ser parafraseada como \textquotedblleft para qualquer moeda, se essa moeda estiver no meu bolso, então ela é uma moeda de vinte e cinco centavos\textquotedblright{}. Assim, podemos simbolizá-la como \textquotedblleft $\forall x(\atom{P}{x} \eif \atom{Q}{x})$\textquotedblright{}.

Como \cref*{quan1} trata de moedas que estão tanto no meu bolso \emph{quanto} são moedas de vinte e cinco centavos, pode ser tentador simbolizá-la usando uma conjunção. No entanto, a sentença \textquotedblleft $\forall x(\atom{P}{x} \eand \atom{Q}{x})$\textquotedblright{} simbolizaria a sentença \textquotedblleft toda moeda é ao mesmo tempo uma moeda de vinte e cinco centavos e está no meu bolso\textquotedblright{}. Isso obviamente significa algo muito diferente de \cref*{quan1}. E assim vemos:
	\factoidbox{
		Uma sentença pode ser simbolizada como $\forall x (\atom{\metav{F}}{x} \eif \atom{\metav{G}}{x})$ se puder ser parafraseada em português como \textquotedblleft todo $F$ é $G$\textquotedblright{}.
	}
\Cref*{quan2} é naturalmente simbolizada usando um quantificador existencial. Ela pode ser parafraseada como \textquotedblleft existe uma moeda que está sobre a mesa e é uma moeda de dez centavos\textquotedblright{}. Assim, podemos simbolizá-la como \textquotedblleft $\exists x(\atom{T}{x} \eand \atom{D}{x})$\textquotedblright{}.

Observe que precisamos usar um condicional com o quantificador universal, mas usamos uma conjunção com o quantificador existencial. Suponha que, em vez disso, tivéssemos escrito \textquotedblleft $\exists x(\atom{T}{x} \eif \atom{D}{x})$\textquotedblright{}. Isso significaria que existe algum objeto no domínio para o qual \textquotedblleft $(\atom{T}{x} \eif \atom{D}{x})$\textquotedblright{} é verdadeira. Lembre-se de que, na LVF, $\metav{A} \eif \metav{B}$ é logicamente equivalente (na LVF) a $\enot\metav{A} \eor \metav{B}$. Essa equivalência também valerá na LPO. Portanto, \textquotedblleft $\exists x(\atom{T}{x} \eif \atom{D}{x})$\textquotedblright{} é verdadeira se houver algum objeto no domínio tal que \textquotedblleft $(\enot \atom{T}{x} \eor \atom{D}{x})$\textquotedblright{} seja verdadeira desse objeto. Isto é, \textquotedblleft $\exists x (\atom{T}{x} \eif \atom{D}{x})$\textquotedblright{} é verdadeira se alguma moeda \emph{ou} não estiver sobre a mesa \emph{ou} for uma moeda de dez centavos. Naturalmente, existe uma moeda que não está sobre a mesa: há moedas em muitos outros lugares. Portanto, é \emph{muito fácil} que \textquotedblleft $\exists x(\atom{T}{x} \eif \atom{D}{x})$\textquotedblright{} seja verdadeira. Um condicional geralmente será o conectivo natural a usar com um quantificador universal, mas um condicional dentro do escopo de um quantificador existencial tende a dizer algo bastante fraco. Como regra geral, não coloque condicionais no escopo de quantificadores existenciais, a menos que tenha certeza de que precisa de um.
	\factoidbox{
		Uma sentença pode ser simbolizada como $\exists x (\atom{\metav{F}}{x} \eand \atom{\metav{G}}{x})$ se puder ser parafraseada em português como \textquotedblleft algum $F$ é~$G$\textquotedblright{}.
	}
\Cref{quan3} pode ser parafraseada como: \textquotedblleft não é o caso de que toda moeda sobre a mesa seja uma moeda de dez centavos\textquotedblright{}. Portanto, podemos simbolizá-la por \textquotedblleft $\enot \forall x(\atom{T}{x} \eif \atom{D}{x})$\textquotedblright{}. Talvez você olhe para \cref*{quan3} e a parafraseie, em vez disso, como \textquotedblleft alguma moeda sobre a mesa não é uma moeda de dez centavos\textquotedblright{}. Então, você a simbolizaria por \textquotedblleft $\exists x(\atom{T}{x} \eand \enot \atom{D}{x})$\textquotedblright{}. Embora provavelmente ainda não seja imediatamente óbvio, essas duas sentenças são logicamente equivalentes. (Isso se deve à equivalência lógica entre $\enot\forall x\,\metav{A}$ e $\exists x\,\enot\metav{A}$, mencionada em \cref{s:FOLBuildingBlocks}, juntamente com a equivalência entre $\enot(\metav{A}\eif\metav{B})$ e $\metav{A}\eand\enot\metav{B}$. )

\Cref{quan4} pode ser parafraseada como: \textquotedblleft não é o caso de que exista alguma moeda de dez centavos no meu bolso\textquotedblright{}. Isso pode ser simbolizado por \textquotedblleft $\enot\exists x(\atom{P}{x} \eand \atom{D}{x})$\textquotedblright{}. Ela também pode ser parafraseada como \textquotedblleft tudo no meu bolso é uma moeda que não é de dez centavos\textquotedblright{} e então poderia ser simbolizada por \textquotedblleft $\forall x(\atom{P}{x} \eif \enot \atom{D}{x})$\textquotedblright{}. Novamente, as duas simbolizações são logicamente equivalentes; ambas são simbolizações corretas de \cref*{quan4}.

\factoidbox{
	Uma sentença que pode ser parafraseada como \textquotedblleft nenhum $F$ é $G$\textquotedblright{} pode ser simbolizada como $\enot\exists x (\atom{\metav{F}}{x} \eand \atom{\metav{G}}{x})$ e também como $\forall x (\atom{\metav{F}}{x} \eif \enot\atom{\metav{G}}{x})$.
}

Por fim, considere \textquotedblleft somente\textquotedblright{}, como em:
\begin{enumerate}
	\item\label{quan5} Somente moedas de dez centavos estão sobre a mesa.
\end{enumerate}
Como devemos simbolizar isso? Uma boa estratégia é considerar quando a sentença seria falsa. Se estamos dizendo que somente moedas de dez centavos estão sobre a mesa, estamos excluindo todos os casos em que algo sobre a mesa não é uma moeda de dez centavos. Assim, podemos simbolizar a sentença da mesma maneira que simbolizaríamos \textquotedblleft nenhuma moeda que não seja de dez centavos está sobre a mesa\textquotedblright{}. Lembrando a lição que acabamos de aprender, e simbolizando \textquotedblleft $x$ não é uma moeda de dez centavos\textquotedblright{} como \textquotedblleft $\enot \atom{D}{x}$\textquotedblright{}, as simbolizações possíveis são: \textquotedblleft $\enot\exists x(\atom{T}{x} \eand \enot \atom{D}{x})$\textquotedblright{} ou, alternativamente, \textquotedblleft $\forall x(\atom{T}{x} \eif \enot\enot \atom{D}{x})$\textquotedblright{}. Como as duplas negações se cancelam, a segunda é tão boa quanto \textquotedblleft $\forall x(\atom{T}{x} \eif \atom{D}{x})$\textquotedblright{}. Em outras palavras, \textquotedblleft somente moedas de dez centavos estão sobre a mesa\textquotedblright{} e \textquotedblleft tudo o que está sobre a mesa é uma moeda de dez centavos\textquotedblright{} são simbolizadas da mesma maneira.

\factoidbox{
	Uma sentença que pode ser parafraseada como \textquotedblleft somente $F$s são $G$s\textquotedblright{} pode ser simbolizada como $\enot\exists x (\atom{\metav{G}}{x} \eand \enot\atom{\metav{F}}{x})$ e também como $\forall x (\atom{\metav{G}}{x} \eif \atom{\metav{F}}{x})$.
}

\section{Predicados vazios}

Em \cref{s:FOLBuildingBlocks}, enfatizamos que um nome deve designar exatamente um objeto no domínio. No entanto, um predicado não precisa aplicar-se a nada no domínio. Um predicado que não se aplica a nenhum objeto no domínio é chamado de \define{predicado vazio}. Vale a pena explorar isso.

\newglossaryentry{empty predicate}{
  name = predicado vazio,
  description = {Um \gls{predicate} que não se aplica a nenhum objeto no \gls{domain}}
}

Suponha que queiramos simbolizar estas duas sentenças:
	\begin{enumerate}
\item\label{monkey1} Todo macaco conhece a linguagem de sinais.
\item\label{monkey2} Algum macaco conhece a linguagem de sinais.
	\end{enumerate}
É possível escrever a chave de simbolização para essas sentenças da seguinte maneira:
	\begin{ekey}
\item[\text{domínio}] animais


	\end{ekey}
\Cref*{monkey1} pode agora ser simbolizada por \textquotedblleft $\forall x(\atom{M}{x} \eif \atom{S}{x})$\textquotedblright{}. \Cref*{monkey2} pode ser simbolizada como \textquotedblleft $\exists x(\atom{M}{x} \eand \atom{S}{x})$\textquotedblright{}.



É tentador dizer que \cref*{monkey1} \emph{implica} \cref*{monkey2}. Isto é, podemos pensar que é impossível que todo macaco conheça a linguagem de sinais sem que algum macaco conheça a linguagem de sinais. Mas isso seria um erro. É possível que a sentença \textquotedblleft $\forall x(\atom{M}{x} \eif \atom{S}{x})$\textquotedblright{} seja verdadeira embora a sentença \textquotedblleft $\exists x(\atom{M}{x} \eand \atom{S}{x})$\textquotedblright{} seja falsa.

Como isso pode acontecer? A resposta vem de considerar se essas sentenças seriam verdadeiras ou falsas \emph{se não houvesse macacos}. Se não houvesse macacos de espécie alguma (no domínio), então \textquotedblleft $\forall x(\atom{M}{x} \eif \atom{S}{x})$\textquotedblright{} seria verdadeira \emph{vacuamente}: escolha qualquer macaco que quiser — ele conhece a linguagem de sinais! Mas, se não houvesse macacos de espécie alguma (no domínio), \textquotedblleft $\exists x(\atom{M}{x} \eand \atom{S}{x})$\textquotedblright{} seria falsa.

Outro exemplo ajudará a esclarecer isso. Suponha que estendamos a chave de simbolização acima, acrescentando:
	\begin{ekey}
\item[\atom{R}{x}] \gap{x} é uma geladeira
	\end{ekey}
Agora considere a sentença \textquotedblleft $\forall x(\atom{R}{x} \eif \atom{M}{x})$\textquotedblright{}. Ela simboliza \textquotedblleft toda geladeira é um macaco\textquotedblright{}. Essa sentença é verdadeira, dada a nossa chave de simbolização, o que é contraintuitivo, pois (presumivelmente) não queremos dizer que haja um monte de macacos-geladeiras. É importante lembrar, porém, que \textquotedblleft $\forall x(\atom{R}{x} \eif \atom{M}{x})$\textquotedblright{} é verdadeira \ifeff{} qualquer membro do domínio que seja uma geladeira seja um macaco. Como o domínio é formado por \emph{animais}, não há geladeiras no domínio. Novamente, então, a sentença é verdadeira \emph{vacuamente}.

Se estivéssemos realmente lidando com a sentença \textquotedblleft todas as geladeiras são macacos\textquotedblright{}, provavelmente desejaríamos incluir eletrodomésticos de cozinha no domínio. Então o predicado \textquotedblleft $R$\textquotedblright{} não seria vazio e a sentença \textquotedblleft $\forall x(\atom{R}{x} \eif \atom{M}{x})$\textquotedblright{} seria falsa.
	\factoidbox{
		Quando $\metav{F}$ é um predicado vazio, qualquer sentença $\forall x (\atom{\metav{F}}{x} \eif \ldots)$ é verdadeira vacuamente.
	}


\section{Escolhendo um domínio}
A simbolização apropriada de uma sentença em português na LPO dependerá da chave de simbolização. Escolher uma chave pode ser difícil. Suponha que queiramos simbolizar a sentença:
	\begin{enumerate}
\item\label{pickdomainrose} Toda rosa tem um espinho.
	\end{enumerate}
Poderíamos oferecer esta chave de simbolização:
	\begin{ekey}
\item[\atom{R}{x}] \gap{x} é uma rosa
\item[\atom{T}{x}] \gap{x} tem um espinho
	\end{ekey}
Pode ser tentador dizer que \cref*{pickdomainrose} deve ser simbolizada como \textquotedblleft $\forall x(\atom{R}{x} \eif \atom{T}{x})$\textquotedblright{}, mas ainda não escolhemos um domínio. Se o domínio contém todas as rosas, essa seria uma boa simbolização. Contudo, se o domínio for apenas \emph{coisas sobre a mesa da minha cozinha}, então \textquotedblleft $\forall x(\atom{R}{x} \eif \atom{T}{x})$\textquotedblright{} chegaria apenas perto de abranger o fato de que toda rosa \emph{sobre a mesa da minha cozinha} tem um espinho. Se não houver rosas sobre a mesa da minha cozinha, a sentença será trivialmente verdadeira. Isso não é o que queremos. Para simbolizar \cref*{pickdomainrose} adequadamente, precisamos incluir todas as rosas no domínio, mas agora temos duas opções.

Primeiro, podemos restringir o domínio para incluir todas as rosas, mas \emph{somente} rosas. Então \cref*{pickdomainrose} pode, se quisermos, ser simbolizada por \textquotedblleft $\forall x\,\atom{T}{x}$\textquotedblright{}. Isso é verdadeiro \ifeff{} tudo no domínio tem um espinho; como o domínio é constituído apenas pelas rosas, isso é verdadeiro \ifeff{} toda rosa tem um espinho. Ao restringir o domínio, conseguimos simbolizar nossa sentença em português com uma sentença muito curta da LPO. Assim, essa abordagem pode nos poupar trabalho, se toda sentença que quisermos tratar for sobre rosas.

Segundo, podemos deixar que o domínio contenha coisas além de rosas: rododendros, ratos, rifles, qualquer coisa; e certamente precisaremos incluir um domínio mais abrangente se quisermos simultaneamente simbolizar sentenças como:
	\begin{enumerate}
\item\label{pickdomaincowboy} Todo vaqueiro canta uma canção triste, triste.
	\end{enumerate}
Nosso domínio agora deve incluir todas as rosas (para que possamos simbolizar \cref*{pickdomainrose}) e todos os vaqueiros (para que possamos simbolizar \cref*{pickdomaincowboy}). Assim, poderíamos oferecer a seguinte chave de simbolização:
	\begin{ekey}
\item[\text{domínio}] pessoas e plantas
\item[\atom{C}{x}] \gap{x} é um vaqueiro
\item[\atom{S}{x}] \gap{x} canta uma canção triste, triste
\item[\atom{R}{x}] \gap{x} é uma rosa
\item[\atom{T}{x}] \gap{x} tem um espinho
	\end{ekey}
Agora teremos de simbolizar \cref*{pickdomainrose} com \textquotedblleft $\forall x (\atom{R}{x} \eif \atom{T}{x})$\textquotedblright{}, pois \textquotedblleft $\forall x\, \atom{T}{x}$\textquotedblright{} simbolizaria a sentença \textquotedblleft toda pessoa ou planta tem um espinho\textquotedblright{}. Da mesma forma, teremos de simbolizar \cref*{pickdomaincowboy} com \textquotedblleft $\forall x (\atom{C}{x} \eif \atom{S}{x})$\textquotedblright{}.

Em geral, o quantificador universal pode ser usado para simbolizar a expressão em português \textquotedblleft todos\textquotedblright{} se o domínio contiver apenas pessoas. Se houver pessoas e outras coisas no domínio, então \textquotedblleft todos\textquotedblright{} deve ser tratado como \textquotedblleft toda pessoa\textquotedblright{}.


\section{A utilidade da paráfrase}
Ao simbolizar sentenças em português na LPO, é importante compreender a estrutura das sentenças que você quer simbolizar. O que importa é a simbolização final na LPO e, às vezes, você poderá passar diretamente de uma sentença em português para uma sentença da LPO. Em outras ocasiões, ajuda parafrasear a sentença uma ou mais vezes. Cada paráfrase sucessiva deve aproximar-se da sentença original de algo que você possa simbolizar diretamente na LPO com facilidade.

Nos próximos exemplos, usaremos esta chave de simbolização:
	\begin{ekey}
\item[\text{domínio}] pessoas
\item[\atom{B}{x}] \gap{x} é baixista
\item[\atom{R}{x}] \gap{x} é uma estrela do rock
		\item[k] Kim Deal
	\end{ekey}
Agora considere estas sentenças:
	\begin{enumerate}
\item\label{pronoun1} Se Kim Deal é baixista, então ela é uma estrela do rock.
\item\label{pronoun2} Se uma pessoa é baixista, então ela é uma estrela do rock.
	\end{enumerate}
As mesmas palavras aparecem como consequente em \cref*{pronoun1,pronoun2} (\textquotedblleft \ldots ela é uma estrela do rock\textquotedblright{}), mas significam coisas muito diferentes. Para deixar isso claro, muitas vezes ajuda parafrasear as sentenças originais, removendo os pronomes.

\Cref*{pronoun1} pode ser parafraseada como: \textquotedblleft Se Kim Deal é baixista, então \emph{Kim Deal} é uma estrela do rock\textquotedblright{}. Isso obviamente pode ser simbolizado como \textquotedblleft $\atom{B}{k} \eif \atom{R}{k}$\textquotedblright{}.

\Cref*{pronoun2} deve ser parafraseada de maneira diferente: \textquotedblleft Se uma pessoa é baixista, então \emph{essa pessoa} é uma estrela do rock\textquotedblright{}. Essa sentença não trata de nenhuma pessoa específica, portanto precisamos de uma variável. Como passo intermediário, podemos parafraseá-la como: \textquotedblleft Para qualquer pessoa $x$, se $x$ é baixista, então $x$ é uma estrela do rock\textquotedblright{}. Agora isso pode ser simbolizado como \textquotedblleft $\forall x (\atom{B}{x} \eif \atom{R}{x})$\textquotedblright{}. Essa é a mesma sentença que usaríamos para simbolizar \textquotedblleft toda pessoa que é baixista é uma estrela do rock\textquotedblright{}. Refletindo, isso certamente é verdadeiro \ifeff{} \cref*{pronoun2} é verdadeira, como esperaríamos.

Considere estas outras sentenças:
	\begin{enumerate}
\item\label{anyone1} Se alguém é baixista, então Kim Deal é uma estrela do rock.
\item\label{anyone2} Se alguém é baixista, então ela é uma estrela do rock.
	\end{enumerate}
As mesmas palavras aparecem como antecedente em \cref*{anyone1,anyone2} (\textquotedblleft se alguém é baixista\ldots\textquotedblright{}), mas pode ser difícil descobrir como simbolizar esses dois usos. Novamente, a paráfrase nos ajudará.

\Cref*{anyone1} pode ser parafraseada como: \textquotedblleft Se há pelo menos um baixista, então Kim Deal é uma estrela do rock\textquotedblright{}. Agora fica claro que este é um condicional cujo antecedente é uma expressão quantificada; portanto, podemos simbolizar a sentença inteira com um condicional como operador lógico principal: \textquotedblleft $\exists x \atom{B}{x} \eif \atom{R}{k}$\textquotedblright{}.

\Cref*{anyone2} pode ser parafraseada como: \textquotedblleft Para todas as pessoas $x$, se $x$ é baixista, então $x$ é uma estrela do rock\textquotedblright{}. Ou, em português mais natural, pode ser parafraseada como \textquotedblleft todos os baixistas são estrelas do rock\textquotedblright{}. É melhor simbolizá-la como \textquotedblleft $\forall x(\atom{B}{x} \eif \atom{R}{x})$\textquotedblright{}, assim como \cref*{pronoun2}.

A lição é que as palavras \textquotedblleft qualquer\textquotedblright{} e \textquotedblleft alguém\textquotedblright{} em português normalmente devem ser simbolizadas usando quantificadores; se você estiver tendo dificuldade para determinar se deve usar um quantificador existencial ou universal, tente parafrasear a sentença com uma frase em português que use palavras \emph{além de} \textquotedblleft qualquer\textquotedblright{} ou \textquotedblleft alguém\textquotedblright{}.



\section{Quantificadores e escopo}
Continuando o exemplo, suponha que queiramos simbolizar estas sentenças:
	\begin{enumerate}
\item\label{qscope1} Se todos são baixistas, então Lars é baixista.
\item\label{qscope2} Todos são tais que, se são baixistas, então Lars é baixista.
	\end{enumerate}
Para simbolizar essas sentenças, teremos de acrescentar um novo nome à chave de simbolização, a saber:
	\begin{ekey}
		\item[l] Lars
	\end{ekey}
\Cref*{qscope1} é um condicional cujo antecedente é \textquotedblleft todos são baixistas\textquotedblright{}, portanto vamos simbolizá-la por \textquotedblleft $\forall x\, \atom{B}{x} \eif \atom{B}{l}$\textquotedblright{}. Essa sentença é \emph{necessariamente} verdadeira: se \emph{todos} são de fato baixistas, escolha qualquer pessoa — por exemplo, Lars — e ela será baixista.

\Cref*{qscope2}, em contraste, pode ser melhor parafraseada como \textquotedblleft toda pessoa $x$ é tal que, se $x$ é baixista, então Lars é baixista\textquotedblright{}.
Isso é simbolizado por \textquotedblleft $\forall x (\atom{B}{x} \eif \atom{B}{l})$\textquotedblright{}.

Essa sentença pode perfeitamente ser falsa. Por exemplo, Kim Deal é baixista.
Assim, \textquotedblleft $\atom{B}{k}$\textquotedblright{} é verdadeira. Suponha que Lars não seja baixista (digamos, que seja baterista em vez disso), de modo que \textquotedblleft $\atom{B}{l}$\textquotedblright{} seja falsa.
Consequentemente, \textquotedblleft $\atom{B}{k} \eif \atom{B}{l}$\textquotedblright{} será falsa, e portanto \textquotedblleft $\forall x
(\atom{B}{x} \eif \atom{B}{l})$\textquotedblright{} também será falsa. Esse exemplo
mostra algo mais: \textquotedblleft $\forall x (\atom{B}{x} \eif \atom{B}{l})$\textquotedblright{} é
falsa exatamente quando $\exists x\,\atom{B}{x} \eif \atom{B}{l}$ é falsa;
elas são equivalentes. Recordando a lição sobre como simbolizar \textquotedblleft qualquer\textquotedblright{}
em \cref{anyone1} da seção anterior, podemos concluir que
\cref*{qscope2} é apenas uma maneira convoluta de dizer \textquotedblleft Lars é baixista,
se alguém o for\textquotedblright{}.


Em resumo, \textquotedblleft $\forall x \atom{B}{x} \eif \atom{B}{l}$\textquotedblright{} e \textquotedblleft $\forall x (\atom{B}{x} \eif \atom{B}{l})$\textquotedblright{} são sentenças muito diferentes. Podemos explicar a diferença em termos do \emph{escopo} do quantificador. O escopo da quantificação é muito semelhante ao escopo da negação, que consideramos ao discutir a LVF, e será útil explicá-lo dessa maneira.

Na sentença \textquotedblleft $\enot \atom{B}{k} \eif \atom{B}{l}$\textquotedblright{}, o escopo de \textquotedblleft $\enot$\textquotedblright{} é apenas o antecedente do condicional. Estamos dizendo algo como: se \textquotedblleft $\atom{B}{k}$\textquotedblright{} é falsa, então \textquotedblleft $\atom{B}{l}$\textquotedblright{} é verdadeira. Da mesma forma, na sentença \textquotedblleft $\forall x \atom{B}{x} \eif \atom{B}{l}$\textquotedblright{}, o escopo de \textquotedblleft $\forall x$\textquotedblright{} é apenas o antecedente do condicional. Estamos dizendo algo como: se \textquotedblleft $\atom{B}{x}$\textquotedblright{} for verdadeira de \emph{tudo}, então \textquotedblleft $\atom{B}{l}$\textquotedblright{} também será verdadeira.

Na sentença \textquotedblleft $\enot(\atom{B}{k} \eif \atom{B}{l})$\textquotedblright{}, o escopo de \textquotedblleft $\enot$\textquotedblright{} é a sentença inteira. Estamos dizendo algo como: \textquotedblleft $(\atom{B}{k} \eif \atom{B}{l})$\textquotedblright{} é falsa. Da mesma forma, na sentença \textquotedblleft $\forall x (\atom{B}{x} \eif \atom{B}{l})$\textquotedblright{}, o escopo de \textquotedblleft $\forall x$\textquotedblright{} é a sentença inteira. Estamos dizendo algo como: \textquotedblleft $(\atom{B}{x} \eif \atom{B}{l})$\textquotedblright{} é verdadeira de \emph{tudo} .

A lição da história é simples. Ao usar condicionais, tenha muito cuidado para certificar-se de que você determinou corretamente o escopo.

\section{Predicados ambíguos}

Suponha que queiramos apenas simbolizar esta sentença:
\begin{enumerate}
\item\label{surgeon1} Adina é uma cirurgiã habilidosa.
\end{enumerate}
Seja o domínio formado por pessoas, seja $\atom{K}{x}$ a expressão que significa \textquotedblleft $x$ é uma cirurgiã habilidosa\textquotedblright{} e seja $a$ um nome para Adina. \Cref{surgeon1} é simplesmente $\atom{K}{a}$.


Suponha, em vez disso, que queiramos simbolizar este argumento:
\begin{earg}
\item O hospital só contratará uma cirurgiã habilidosa.
\item Todas as cirurgiãs são gananciosas.
\item Billy é cirurgião, mas não é habilidoso.
\item[\texttherefore] Billy é ganancioso, mas o hospital não o contratará.
\end{earg}
Precisamos distinguir entre ser uma \emph{cirurgiã habilidosa} e ser simplesmente uma \emph{cirurgiã}. Portanto, definimos esta chave de simbolização:
\begin{ekey}
\item[\text{domínio}] pessoas
\item[\atom{G}{x}] \gap{x} é ganancioso
\item[\atom{H}{x}] O hospital contratará \gap{x}
\item[\atom{R}{x}] \gap{x} é cirurgião
\item[\atom{K}{x}] \gap{x} é habilidoso
\item[b] Billy
\end{ekey}

Agora o argumento pode ser simbolizado da seguinte maneira:
\begin{earg}
\label{surgeon2}
\item[] $\forall x\bigl[\enot (\atom{R}{x} \eand \atom{K}{x}) \eif \enot \atom{H}{x}\bigr]$
\item[] $\forall x(\atom{R}{x} \eif \atom{G}{x})$
\item[] $\atom{R}{b} \eand \enot \atom{K}{b}$
\item[\texttherefore] $\atom{G}{b} \eand \enot \atom{H}{b}$
\end{earg}

Em seguida, suponha que queiramos simbolizar este argumento:
\begin{earg}
\label{surgeon3}
\item Carol é uma cirurgiã habilidosa e uma jogadora de tênis.
\item[\texttherefore] Carol é uma jogadora de tênis habilidosa.
\end{earg}
Se começarmos com a chave de simbolização usada para o argumento anterior, poderemos acrescentar um predicado (seja $\atom{T}{x}$ a expressão que significa \textquotedblleft $x$ é jogador de tênis\textquotedblright{}) e um nome (seja \textquotedblleft $c$\textquotedblright{} um nome para Carol). Então, o argumento se torna:
\begin{earg}
\item[] $(\atom{R}{c} \eand \atom{K}{c}) \eand \atom{T}{c}$
\item[\texttherefore] $\atom{T}{c} \eand \atom{K}{c}$
\end{earg}
Essa simbolização é um desastre! Ela toma um argumento terrível em inglês e o simboliza como um argumento válido na LPO. O problema é que há uma diferença entre ser \emph{habilidoso como cirurgião} e ser \emph{habilidoso como jogador de tênis}. A simbolização correta desse argumento exige dois predicados distintos, um para cada tipo de habilidade. Se fizermos $\atom{K_1}{x}$ significar \textquotedblleft $x$ é habilidoso como cirurgião\textquotedblright{} e $\atom{K_2}{x}$ significar \textquotedblleft $x$ é habilidoso como jogador de tênis\textquotedblright{}, poderemos simbolizar o argumento da seguinte maneira:
\begin{earg}
\label{surgeon3correct}
\item[] $(\atom{R}{c} \eand \atom{K_1}{c}) \eand \atom{T}{c}$
\item[\texttherefore] $\atom{T}{c} \eand \atom{K_2}{c}$
\end{earg}
Assim como o argumento em inglês que simboliza, este é inválido. %\nix{Observe que não há conexão lógica entre $\atom{K_1}{c}$ e $\atom{R}{c}$. Como símbolos da LPO, eles poderiam ser quaisquer predicados de um lugar. Em inglês, há uma conexão entre ser um \emph{cirurgião} e ser uma \emph{cirurgiã habilidosa}: toda cirurgiã habilidosa é cirurgiã. Para capturar essa conexão, simbolizamos \textquotedblleft Carol é uma cirurgiã habilidosa\textquotedblright{} como $\atom{R}{c} \eand K_1c$. Isso significa: \textquotedblleft Carol é cirurgiã e é habilidosa como cirurgiã.\textquotedblright{}}

A lição desses exemplos é que precisamos ter cuidado ao simbolizar predicados de maneira ambígua. Problemas semelhantes podem surgir com predicados como \emph{bom}, \emph{mau}, \emph{grande} e \emph{pequeno}. Assim como cirurgiões habilidosos e jogadores de tênis habilidosos têm habilidades diferentes, cães grandes, ratos grandes e problemas grandes são grandes de maneiras diferentes.

É suficiente ter um predicado que signifique \textquotedblleft $x$ é uma cirurgiã habilidosa\textquotedblright{}, em vez de dois predicados \textquotedblleft $x$ é habilidoso\textquotedblright{} e \textquotedblleft $x$ é cirurgião\textquotedblright{}? Às vezes. Como \cref{surgeon1} mostra, às vezes não precisamos distinguir entre cirurgiões habilidosos e outros cirurgiões.

Devemos sempre distinguir entre diferentes maneiras de ser habilidoso, bom, mau ou grande? Não. Como mostra o argumento sobre Billy, às vezes precisamos falar apenas de um tipo de habilidade. Se você está simbolizando um argumento que trata apenas de cães, não há problema em definir um predicado que signifique \textquotedblleft $x$ é grande\textquotedblright{}. Se o domínio inclui cães e ratos, porém, provavelmente é melhor fazer o predicado significar \textquotedblleft $x$ é grande para um cão\textquotedblright{}.

\practiceproblems
\problempart
\phantomsection\label{pr.BarbaraEtc}
Eis as figuras silogísticas identificadas por Aristóteles e seus sucessores, juntamente com seus nomes medievais:
\begin{compactlist}
	\item \textbf{Barbara.} Todo $G$ é $F$. Todo $H$ é $G$. Portanto: todo $H$ é $F$.
	\item \textbf{Celarent.} Nenhum $G$ é $F$. Todo $H$ é $G$. Portanto: nenhum $H$ é $F$.
	\item \textbf{Ferio.} Nenhum $G$ é $F$. Algum $H$ é $G$. Portanto: algum $H$ não é $F$.
	\item \textbf{Darii.} Todo $G$ é $F$. Algum $H$ é $G$. Portanto: algum $H$ é $F$.
	\item \textbf{Camestres.} Todo $F$ é $G$. Nenhum $H$ é $G$. Portanto: nenhum $H$ é $F$.
	\item \textbf{Cesare.} Nenhum $F$ é $G$. Todo $H$ é $G$. Portanto: nenhum $H$ é $F$.
	\item \textbf{Baroko.} Todo $F$ é $G$. Algum $H$ não é $G$. Portanto: algum $H$ não é~$F$.
	\item \textbf{Festino.} Nenhum $F$ é $G$. Algum $H$ é $G$. Portanto: algum $H$ não é~$F$.
	\item \textbf{Datisi.} Todo $G$ é $F$. Algum $G$ é $H$. Portanto: algum $H$ é $F$.
	\item \textbf{Disamis.} Algum $G$ é $F$. Todo $G$ é $H$. Portanto: algum $H$ é $F$.
	\item \textbf{Ferison.} Nenhum $G$ é $F$. Algum $G$ é $H$. Portanto: algum $H$ não é $F$.
	\item \textbf{Bokardo.} Algum $G$ não é $F$. Todo $G$ é $H$. Portanto: algum $H$ não é~$F$.
	\item \textbf{Camenes.} Todo $F$ é $G$. Nenhum $G$ é $H$. Portanto: nenhum $H$ é $F$.
	\item \textbf{Dimaris.} Algum $F$ é $G$. Todo $G$ é $H$. Portanto: algum $H$ é $F$.
	\item \textbf{Fresison.} Nenhum $F$ é $G$. Algum $G$ é $H$. Portanto: algum $H$ não é~$F$.
\end{compactlist}
Simbolize cada figura na LPO.

\

\problempart
\label{pr.FOLvegetarians}
Usando a seguinte chave de simbolização:
\begin{ekey}
\item[\text{domínio}] pessoas
\item[\atom{K}{x}] \gap{x} conhece a combinação do cofre
\item[\atom{S}{x}] \gap{x} é espião
\item[\atom{V}{x}] \gap{x} é vegetariano
%\item[\atom{T}{x,y}] \gap{x} confia em \gap{y}.
\item[h] Hofthor
\item[i] Ingmar
\end{ekey}
simbolize as seguintes sentenças na LPO:
\begin{compactlist}
\item Nem Hofthor nem Ingmar é vegetariano.
\item Nenhum espião conhece a combinação do cofre.
\item Ninguém conhece a combinação do cofre a menos que Ingmar a conheça.
\item Hofthor é espião, mas nenhum vegetariano é espião.
\end{compactlist}
\solutions
\problempart\label{pr.FOLalligators}
Usando esta chave de simbolização:
\begin{ekey}
\item[\text{domínio}] todos os animais
\item[\atom{A}{x}] \gap{x} é um jacaré
\item[\atom{M}{x}] \gap{x} é um macaco
\item[\atom{R}{x}] \gap{x} é um réptil
\item[\atom{Z}{x}] \gap{x} vive no zoológico
\item[a] Amos
\item[b] Bouncer
\item[c] Cleo
\end{ekey}
simbolize cada uma das seguintes sentenças na LPO:
\begin{compactlist}
\item Amos, Bouncer e Cleo vivem no zoológico.
\item Bouncer é um réptil, mas não é um jacaré.
% \item Se Cleo ama Bouncer, então Bouncer é um macaco.
% \item Se Bouncer e Cleo são jacarés, então Amos ama ambos.
\item Algum réptil vive no zoológico.
\item Todo jacaré é um réptil.
\item Todo animal que vive no zoológico é um macaco ou um jacaré.
\item Há répteis que não são jacarés.
% \item Cleo ama um réptil.
% \item Bouncer ama todos os macacos que vivem no zoológico.
% \item Todos os macacos que Amos ama também o amam.
\item Se algum animal é um réptil, então Amos é.
\item Se algum animal é um jacaré, então ele é um réptil.
% \item Todo macaco que Cleo ama também é amado por Amos.
% \item Há um macaco que ama Bouncer, mas infelizmente Bouncer não retribui esse amor.
\end{compactlist}

\problempart
\label{pr.FOLarguments}
Para cada argumento, escreva uma chave de simbolização e simbolize o argumento na LPO.
\begin{compactlist}
\item Willard é lógico. Todos os lógicos usam chapéus engraçados. Portanto, Willard usa um chapéu engraçado.
\item Nada sobre a minha mesa escapa à minha atenção. Há um computador sobre a minha mesa. Portanto, há um computador que não escapa à minha atenção.
\item Todos os meus sonhos são em preto e branco. Programas de televisão antigos são em preto e branco. Portanto, alguns dos meus sonhos são programas de televisão antigos.
\item Nem Holmes nem Watson estiveram na Austrália. Uma pessoa só poderia ter visto um canguru se tivesse estado na Austrália ou em um zoológico. Embora Watson não tenha visto um canguru, Holmes viu. Portanto, Holmes esteve em um zoológico.
\item Ninguém espera a Inquisição Espanhola. Ninguém conhece os problemas que enfrentei. Portanto, qualquer pessoa que espera a Inquisição Espanhola conhece os problemas que enfrentei.
\item Todos os bebês são ilógicos. Ninguém que seja ilógico consegue lidar com um crocodilo. Berthold é um bebê. Portanto, Berthold não consegue lidar com um crocodilo.
\end{compactlist}

% \problempart
% label{pr.QLarguments}
% Para cada argumento, escreva uma chave de simbolização e simbolize o argumento na LPO.
% \begin{compactlist}
% \item Nada sobre a minha mesa escapa à minha atenção. Há um computador sobre a minha mesa. Portanto, há um computador que não escapa à minha atenção.
% \item Todos os meus sonhos são em preto e branco. Programas de televisão antigos são em preto e branco. Portanto, alguns dos meus sonhos são programas de televisão antigos.
% \item Nem Holmes nem Watson estiveram na Austrália. Uma pessoa só poderia ver um canguru se tivesse estado na Austrália ou em um zoológico. Embora Watson não tenha visto um canguru, Holmes viu. Portanto, Holmes esteve em um zoológico.
% \item Ninguém espera a Inquisição Espanhola. Ninguém conhece os problemas que enfrentei. Portanto, qualquer pessoa que espera a Inquisição Espanhola conhece os problemas que enfrentei.
% \item Um antílope é maior que uma caixa de pão. Estou pensando em algo que não é maior que uma caixa de pão, e é um antílope ou um melão. Portanto, estou pensando em um melão.
% \item Todos os bebês são ilógicos. Ninguém que seja ilógico consegue lidar com um crocodilo. Berthold é um bebê. Portanto, Berthold não consegue lidar com um crocodilo.
% \end{compactlist}

\chapter{Multiple generality}\label{s:MultipleGenerality}
So far, we have only considered sentences that require one-place predicates and one quantifier. The full power of FOL really comes out when we start to use many-place predicates and multiple quantifiers. For this insight, we largely have \foreignlanguage{german}{Gottlob Frege} (1879) to thank, but also C.~S.~Peirce.


\section{Many-placed predicates}
All of the predicates that we have considered so far concern properties that objects might have. Those predicates have one gap in them, and to make a sentence, we simply need to slot in one term. They are \define{one-place} predicates.

However, other predicates concern the \emph{relation} between two things. Here are some examples of relational predicates in English:
	\begin{itemize}
		\item \blank\ loves \blank
		\item \blank\ is to the left of \blank
		\item \blank\ is in debt to \blank
	\end{itemize}
These are \define{two-place} predicates. They need to be filled in with two terms in order to make a sentence. Conversely, if we start with an English sentence containing many singular terms, we can remove two singular terms, to obtain different two-place predicates. Consider the sentence `Vinnie borrowed the family car from Nunzio'. By deleting two singular terms, we can obtain any of three different two-place predicates
	\begin{itemize}
		\item Vinnie borrowed \blank\ from \blank
		\item \blank\ borrowed the family car from \blank
		\item \blank\ borrowed \blank\ from Nunzio
	\end{itemize}
and by removing all three singular terms, we  obtain a \define{three-place} predicate:
	\begin{itemize}
		\item \blank\ borrowed \blank\ from \blank
	\end{itemize}
Indeed, there is no in principle upper limit on the number of places that our predicates may contain.

\section{Mind the gap(s)!}

We have used the same symbol, `\blank', to indicate a gap formed by deleting a term from a sentence. However, as \foreignlanguage{german}{Frege} emphasized, these are \emph{different} gaps. To obtain a sentence, we can fill them in with the same term, but we can equally fill them in with different terms, and in various different orders. The following are three perfectly good sentences, obtained by filling in the gaps in `\blank\ loves \blank{}' in different ways; but they all have distinctively different meanings:
\begin{enumerate}
	\item\label{terms3} Karl loves Imre.
	\item\label{terms3b} Imre loves Karl.
	\item\label{terms3a} Karl loves Karl.
\end{enumerate}
The point is that we need to keep track of the gaps in predicates, so that we can keep track of how we are filling them in. To keep track of the gaps, we assign them variables. Suppose we want to symbolize the preceding sentences. Then I might start with the following symbolization key: 
	\begin{ekey}
		\item[\text{domain}] people
		\item[i] Imre
		\item[k] Karl
		\item[\atom{L}{x,y}] \gap{x} loves \gap{y}
	\end{ekey}
 \Cref*{terms3} will be symbolized by `$\atom{L}{k,i}$', \cref*{terms3b} will be symbolized by `$\atom{L}{i,k}$', and \cref*{terms3a} will be symbolized by `$\atom{L}{k,k}$'. Here are a few more sentences that we can symbolize with the same key:
\begin{enumerate}
	\item\label{terms4} Imre loves himself.
	\item\label{terms5} Karl loves Imre, but not vice versa.
	\item\label{terms6} Karl is loved by Imre.
\end{enumerate}
\Cref*{terms4} can be paraphrased as `Imre loves Imre', and so
symbolized by `$\atom{L}{i,i}$'. \Cref*{terms5} is a conjunction. We
can paraphrase it as `Karl loves Imre, and Imre does not love Karl',
and so symbolize it as `$\atom{L}{k,i} \eand \enot \atom{L}{i,k}$'.
\Cref*{terms6} can be paraphrased by `Imre loves Karl', and so
symbolized as `$\atom{L}{i,k}$'. In this last case, of course, we have
lost the difference in \emph{tone} between the active and passive
voice; but we have at least preserved the truth conditions.

But the relationship between `Imre loves Karl' and `Karl is loved by Imre' highlights something important. To see what, suppose we add another entry to our symbolization key:
\begin{ekey}
	\item[\atom{M}{x,y}] \gap{y} loves \gap{x}
\end{ekey}
The entry for `$M$' uses exactly the same English word---`loves'---as the entry for `$L$'. \emph{But the gaps have been swapped around!} (Just look closely at the subscripts.) And this \emph{matters}.

To explain: when we see a sentence like `$\atom{L}{k,i}$', we are
being told to take the \emph{first} name (i.e., `$k$') and associate
its value (i.e., Karl) with the gap labelled `$x$', then take the
\emph{second} name (i.e.,~`$i$') and associate its value (i.e., Imre)
with the gap labelled `$y$', and so come up with: \emph{Karl loves
Imre}. The sentence `$\atom{M}{i,k}$' also tells us to take the
\emph{first} name (i.e., `$i$') and plug its value into the gap
labelled `$x$', and take the \emph{second} name (i.e., `$k$') and plug
its value into the gap labelled `$y$', and so also come up with:
\emph{Karl loves Imre}.

So, `$\atom{L}{i,k}$' and `$\atom{M}{k,i}$' both symbolize `Imre loves Karl', whereas `$\atom{L}{k,i}$' and `$\atom{M}{i,k}$' both symbolize `Karl loves Imre'. Since love can be unrequited, these are different claims.

One last example might be helpful. Suppose we add this to our symbolisation key: 
\begin{ekey}
	\item[\atom{P}{x,y}] \gap{x} prefers \gap{x} to \gap{y}
\end{ekey}
Now the sentence `$\atom{P}{i,k}$' symbolizes `Imre prefers Imre to Karl', and `$\atom{P}{k,i}$' symbolizes `Karl prefers Karl to Imre'.  And note that we could have achieved the same effect, if we had instead specified:
\begin{ekey}
	\item[\atom{P}{x,y}] \gap{x} prefers themselves to \gap{y}
\end{ekey}
In any case, the overall moral of this is simple. \emph{When dealing with predicates with more than one place, pay careful attention to the order of the gaps!} 


\section{The order of quantifiers}\label{ss:OrderQuant}
Consider the sentence `everyone loves someone'. This is potentially ambiguous. It might mean either of the following:
	\begin{enumerate}
		\item\label{lovecycle} For every person x, there is some person that x loves.
		\item\label{loveconverge} There is some particular person whom every person loves.
	\end{enumerate}
\Cref*{lovecycle} can be symbolized by `$\forall x \exists y\, \atom{L}{x,y}$'. Suppose that our domain of discourse is restricted to Imre, Juan and Karl. Suppose also that Karl loves Imre but not Juan, that Imre loves Juan but not Karl, and that Juan loves Karl but not Imre. Then \cref*{lovecycle} is true.

\Cref*{loveconverge} is symbolized by `$\exists y \forall x\,
\atom{L}{x,y}$'. \Cref{loveconverge} is \emph{not} true in the
situation just described, since no single person is loved by everyone.
All three of Juan, Imre and Karl would have to converge on (at least)
one object of love.

The point of the example is to illustrate that the order of the quantifiers matters a great deal. Indeed, to switch them around is called a \emph{quantifier shift fallacy}. Here is an example, which comes up in various forms throughout the philosophical literature:
	\begin{earg}
		\item[] For every person, there is some truth they cannot know. \hfill ($\forall \exists$)
		\item[\texttherefore] There is some particular truth that no person can know. \hfill ($\exists \forall$)
	\end{earg}
This argument form is obviously invalid. It's just as bad as:\footnote{Thanks to Rob Trueman for the example.}
	\begin{earg}
		\item[] Every dog has its day. \hfill ($\forall \exists$)
		\item[\texttherefore] There is a day for all the dogs. \hfill ($\exists \forall$)
	\end{earg}

The order of quantifiers is also important in definitions in mathematics.  For instance, there is a big difference between pointwise and uniform continuity of functions:
\begin{itemize}
\item A function $f$ is \emph{pointwise continuous} if
\[
\forall \epsilon\forall x\exists \delta\forall y(\left|x - y\right| < \delta \to \left|f(x) - f(y)\right| < \epsilon)
\]
\item A function $f$ is \emph{uniformly continuous} if
\[
\forall \epsilon\exists \delta\forall x\forall y(\left|x - y\right| < \delta \to \left|f(x) - f(y)\right| < \epsilon)
\]
\end{itemize}

The moral is simple: \emph{take great care with the order of your quantifiers!}


\section{Stepping-stones to symbolization}
As we are starting to see, symbolization in FOL can become a bit tricky. So, when symbolizing a complex sentence, you should lay down several stepping-stones. As usual, the idea is best illustrated by example. Consider this symbolisation key: 
\begin{ekey}
\item[\text{domain}] people and dogs
\item[\atom{D}{x}] \gap{x} is a dog
\item[\atom{F}{x,y}] \gap{x} is a friend of \gap{y}
\item[\atom{O}{x,y}] \gap{x} owns \gap{y}
\item[g] Geraldo
\end{ekey}
Now let's try to symbolize these sentences:
\begin{enumerate}
\item\label{dog2} Geraldo is a dog owner.
\item\label{dog3} Someone is a dog owner.
\item\label{dog4} All of Geraldo's friends are dog owners.
\item\label{dog5} Every dog owner is a friend of a dog owner.
\item\label{dog6} Every dog owner's friend owns a dog of a friend.
\end{enumerate}
\Cref*{dog2} can be paraphrased as, `There is a dog that Geraldo owns'. This can be symbolized by `$\exists x(\atom{D}{x} \eand \atom{O}{g,x})$'.

\Cref*{dog3} can be paraphrased as, `There is some $y$ such that $y$ is a dog owner'. Dealing with part of this, we might write `$\exists y(y\text{ is a dog owner})$'. Now the fragment we have left as `$y$ is a dog owner' is much like \cref*{dog2}, except that it is not specifically about Geraldo. So we can symbolize \cref*{dog3} by:
$$\exists y \exists x(\atom{D}{x} \eand \atom{O}{y,x})$$
We should pause to clarify something here. In working out how to
symbolize the last sentence, we wrote down `$\exists y(y\text{ is a
dog owner})$'. To be very clear: this is \emph{neither} an FOL
sentence \emph{nor} an English sentence: it uses bits of FOL
(`$\exists$', `$y$') and bits of English (`dog owner'). It is really
\emph{just a stepping-stone} on the way to symbolizing the entire
English sentence with a sentence of FOL. You should regard it as a bit of
rough-working-out, on a par with the doodles that you might
absent-mindedly draw in the margin of this book, whilst you are
concentrating fiercely on some problem.  

\Cref*{dog4} can be paraphrased as, `Every $x$ who is a friend of Geraldo is a dog owner'. Using our stepping-stone tactic, we might write 
$$\forall x \bigl[\atom{F}{x,g} \eif x \text{ is a dog owner}\bigr]$$
Now the fragment that we have left to deal with, `$x$ is a dog owner', is structurally just like \cref*{dog2}. However, it would be a mistake for us simply to write 
$$\forall x \bigl[\atom{F}{x,g} \eif \exists x(\atom{D}{x} \eand \atom{O}{x,x})\bigr]$$
for we would here have a \emph{clash of variables}. The scope of the universal quantifier, `$\forall x$', is the entire conditional, so the `$x$' in `$\atom{D}{x}$' should be governed by that, but `$\atom{D}{x}$' also falls under the scope of the existential quantifier `$\exists x$', so the `$x$' in `$\atom{D}{x}$' should be governed by that. Now confusion reigns: which `$x$' are we talking about? Suddenly the sentence becomes ambiguous (if it is even meaningful at all), and logicians hate ambiguity. The broad moral is that a single variable cannot serve two quantifier-masters simultaneously.

To continue our symbolization, then, we must choose some different variable for our existential quantifier. What we want is something like:
$$\forall x\bigl[\atom{F}{x,g} \eif\exists z(\atom{D}{z} \eand \atom{O}{x,z})\bigr]$$
This adequately symbolizes \cref*{dog4}.

\Cref{dog5} can be paraphrased as `For any $x$ that is a dog owner, there is a dog owner who $x$ is a friend of'. Using our stepping-stone tactic, this becomes 
$$\forall x\bigl[\mbox{$x$ is a dog owner}\eif\exists y(\mbox{$y$ is a dog owner}\eand \atom{F}{x,y})\bigr]$$
Completing the symbolization, we end up with
$$\forall x\bigl[\exists z(\atom{D}{z} \eand \atom{O}{x,z})\eif\exists y\bigl(\exists z(\atom{D}{z} \eand \atom{O}{y,z})\eand \atom{F}{x,y}\bigr)\bigr]$$
Note that we have used the same letter, `$z$', in both the antecedent and the consequent of the conditional, but that these are governed by two different quantifiers. This is ok: there is no clash here, because it is clear which quantifier that variable falls under. We might graphically represent the scope of the quantifiers thus:
$$\overbrace{\forall x\bigl[\overbrace{\exists z(\atom{D}{z} \eand \atom{O}{x,z})}^{\text{scope of 1st `}\exists z\text{'}}\eif \overbrace{\exists y(\overbrace{\exists z(\atom{D}{z} \eand \atom{O}{y,z})}^{\text{scope of 2nd `}\exists z\text{'}}\eand \atom{F}{x,y})\bigr]}^{\text{scope of `}\exists y\text{'}}}^{\text{scope of `}\forall x\text{'}}$$
This shows that no variable is being forced to serve two masters simultaneously.

\Cref{dog6} is the trickiest yet. First we paraphrase it as
`For any $x$ that is a friend of a dog owner, $x$ owns a dog which is
also owned by a friend of~$x$'. Using our stepping-stone tactic, this
becomes:
\ifHTMLtarget
\[\forall x\bigl[x\text{ is a friend of a dog owner}\eif 
x\text{ owns a dog which is owned by a friend of }x\bigr]\]
\else
\begin{multline*}
\forall x\bigl[x\text{ is a friend of a dog owner}\eif {}\\
x\text{ owns a dog which is owned by a friend of }x\bigr]
\end{multline*}
\fi
Breaking this down a bit more:
\ifHTMLtarget
\[\forall x\bigl[\exists y(\atom{F}{x,y} \eand y\text{ is a dog owner})\eif
\exists y((\atom{D}{y} \eand \atom{O}{x,y}) \eand y\text{ is owned by
a friend of }x)\bigr]\]
\else
\begin{multline*}
	\forall x\bigl[\exists y(\atom{F}{x,y} \eand y\text{ is a dog owner})\eif {}\\
\exists y((\atom{D}{y} \eand \atom{O}{x,y}) \eand y\text{ is owned by a friend of }x)\bigr]
\end{multline*}
\fi
And a bit more: 
\ifHTMLtarget
\[\forall x\bigl[\exists y(\atom{F}{x,y} \eand \exists z(\atom{D}{z} \eand \atom{O}{y,z})) \eif 
\exists y((\atom{D}{y} \eand \atom{O}{x,y}) \eand \exists z(\atom{F}{z,x} \eand \atom{O}{z,y}))\bigr]\]
\else
\begin{multline*}
\forall x\bigl[\exists y(\atom{F}{x,y} \eand \exists z(\atom{D}{z} \eand \atom{O}{y,z})) \eif {}\\
\exists y((\atom{D}{y} \eand \atom{O}{x,y}) \eand \exists z(\atom{F}{z,x} \eand \atom{O}{z,y}))\bigr]
\end{multline*}
\fi
And we are done!

There is one subtle issue we should briefly address. We paraphrased
\cref{dog3} as `There is some $y$ such that $y$ is a dog owner'.
Now our domain includes people and dogs, and `someone' includes
people, but (at least arguably) does not include dogs.  To be more
correct, we should have paraphrased \cref*{dog3} as `There is
some $y$ such that $y$ is a person and a dog owner'. A more accurate
symbolization of `Someone is a dog owner' would require that we add a
predicate for `\blank{} is a person' to our symbolization key:
\begin{ekey}
	\item[\atom{P}{x}] \gap{x} is a person
\end{ekey}
Then we can give a better symbolization of \cref{dog3}:
$$\exists y (\atom{P}{y} \eand \exists x(\atom{D}{x} \eand
\atom{O}{y,x}))$$
`Everyone' and `no one' have to be treated similarly: `Everyone
is a friend of Geraldo' and `No one is a friend of Geraldo' would be
symbolized, respectively, as
\begin{align*}
	& \forall x(\atom{P}{x} \eif \atom{F}{x,g})\\
	& \forall x(\atom{P}{x} \eif \enot\atom{F}{x,g}).
\end{align*}
Only `someone', `everyone', and `no one' require this treatment. In
particular, we do not need to explicitly state in the symbolization of
\cref{dog2} that Geraldo is a person. Neither do we have to
ensure in \cref{dog4} that Geraldo's friend~$x$ is a person.
Although it may be true that only people can own dogs or be friends
with Geraldo, it is not part of what the sentences \emph{say}, and so
does not need to be taken into account when we symbolize
them.\footnote{You might object: but the sentences also don't say that
dogs aren't people. E.g., if we're talking about the fictional world
of Mickey Mouse, Goofy should be included in `everyone', but Pluto
should not be.  That's why we picked the predicate `\blank{} is a
person' and not `\blank{} is a human': in that domain, Goofy would
fall under both `\blank{} is a person' and `\blank{} is a dog', but
Pluto would only fall under `\blank{} is a dog'.}

\section{Supressed quantifiers}\label{ss:SuppQuant}

Logic can often help to get clear on the meanings of English claims,
especially where the quantifiers are left implicit or their order is
ambiguous or unclear. The clarity of expression and thinking afforded
by FOL can give you a significant advantage in argument, as can be
seen in the following takedown by British political philosopher Mary
Astell (1666--1731) of her contemporary, the theologian William
Nicholls. In Discourse IV: The Duty of Wives to their Husbands of his
\textit{The Duty of Inferiors towards their Superiors, in Five
  Practical Discourses} (London 1701), Nicholls argued that women are
naturally inferior to men. In the preface to the 3rd edition of her
treatise \textit{Some Reflections upon Marriage, Occasion'd by the Duke
  and Duchess of Mazarine's Case; which is also considered,} Astell
responded as follows:
\begin{quotation}
'Tis true, thro' Want of Learning, and of that Superior Genius which
Men as Men lay claim to, she [Astell] was ignorant of the
\textit{Natural Inferiority} of our Sex, which our Masters lay down as
a Self-Evident and Fundamental Truth. She saw nothing in the Reason of
Things, to make this either a Principle or a Conclusion, but much to
the contrary; it being Sedition at least, if not Treason to assert it
in this Reign.

For if by the Natural Superiority of their Sex, they mean that
\textit{every} Man is by Nature superior to \textit{every} Woman,
which is the obvious meaning, and that which must be stuck to if they
would speak Sense, it wou'd be a Sin in \textit{any} Woman to have
Dominion over \textit{any} Man, and the greatest Queen ought not to
command but to obey her Footman, because no Municipal Laws can
supersede or change the Law of Nature; so that if the Dominion of the
Men be such, the \textit{Salique Law,}\footnote{The Salique law was
  the common law of France which prohibited the crown be passed on to
  female heirs.} as unjust as \textit{English Men} have ever thought
it, ought to take place over all the Earth, and the most glorious
Reigns in the \textit{English, Danish, Castilian}, and other Annals,
were wicked Violations of the Law of Nature!

If they mean that \textit{some} Men are superior to \textit{some}
Women this is no great Discovery; had they turn'd the Tables they
might have seen that \textit{some} Women are Superior to \textit{some}
Men. Or had they been pleased to remember their Oaths of Allegiance
and Supremacy, they might have known that \textit{One} Woman is
superior to \textit{All} the Men in these Nations, or else they have
sworn to very little purpose.\footnote{In 1706, England was ruled by
  Queen Anne.} And it must not be suppos'd, that their Reason and
Religion wou'd suffer them to take Oaths, contrary to the Laws of
Nature and Reason of things.\footnote{Mary Astell, \textit{Reflections
    upon Marriage}, 1706 Preface, iii--iv, and Mary Astell,
  \textit{Political Writings}, Patricia Springborg (ed.), Cambridge
  University Press, 1996, pp.~9--10.}
\end{quotation}
We can symbolize the different interpretations Astell offers of
Nicholls' claim that men are superior to women:
He either meant that every man is superior to every woman, i.e.,
\[
\forall x(\atom{M}{x} \eif \forall y(\atom{W}{y} \eif \atom{S}{x,y}))
\]
or that some men are superior to some women,
\[
\exists x(\atom{M}{x} \eand \exists y(\atom{W}{y} \eand \atom{S}{x,y})).
\]
The latter is true, but so is
\[
\exists y(\atom{W}{y} \eand \exists x(\atom{M}{x} \eand \atom{S}{y,x}))
\]
(some women are superior to some men), so that would be ``no great
discovery.''  In fact, since the Queen is superior to all her
subjects, it's even true that some woman is superior to every man,
i.e.,
\[
\exists y(\atom{W}{y} \land \forall x(\atom{M}{x} \eif \atom{S}{y,x})).
\]
But this is incompatible with the ``obvious meaning'' of Nicholls'
claim, i.e., the first reading. So what Nicholls claims amounts to
treason against the Queen!

\practiceproblems
\solutions
\problempart
Using this symbolization key:
\begin{ekey}
\item[\text{domain}] all animals
\item[\atom{A}{x}] \gap{x} is an alligator
\item[\atom{M}{x}] \gap{x} is a monkey
\item[\atom{R}{x}] \gap{x} is a reptile
\item[\atom{Z}{x}] \gap{x} lives at the zoo
\item[\atom{L}{x,y}] \gap{x} loves \gap{y}
\item[a] Amos
\item[b] Bouncer
\item[c] Cleo
\end{ekey}
symbolize each of the following sentences in FOL:
\begin{compactlist}
%\item Amos, Bouncer, and Cleo all live at the zoo.
%\item Bouncer is a reptile, but not an alligator.
\item If Cleo loves Bouncer, then Bouncer is a monkey.
\item If both Bouncer and Cleo are alligators, then Amos loves them both.
%\item Some reptile lives at the zoo.
%\item Every alligator is a reptile.
%\item Any animal that lives at the zoo is either a monkey or an alligator.
%\item There are reptiles which are not alligators.
\item Cleo loves a reptile.
\item Bouncer loves all the monkeys that live at the zoo.
\item All the monkeys that Amos loves love him back.
%\item If any animal is an reptile, then Amos is.
%\item If any animal is an alligator, then it is a reptile.
\item Every monkey that Cleo loves is also loved by Amos.
\item There is a monkey that loves Bouncer, but sadly Bouncer does not reciprocate this love.
\end{compactlist}

\problempart 
Using the following symbolization key:
\begin{ekey}
\item[\text{domain}] all animals
\item[\atom{D}{x}] \gap{x} is a dog
\item[\atom{S}{x}] \gap{x} likes samurai movies
\item[\atom{L}{x,y}] \gap{x} is larger than \gap{y}
\item[r] Rave
\item[h] Shane
\item[d] Daisy
\end{ekey}
symbolize the following sentences in FOL:
\begin{compactlist}
\item Rave is a dog who likes samurai movies.
\item Rave, Shane, and Daisy are all dogs.
\item Shane is larger than Rave, and Daisy is larger than Shane.
\item All dogs like samurai movies.
\item Only dogs like samurai movies.
\item There is a dog that is larger than Shane.
\item If there is a dog larger than Daisy, then there is a dog larger than Shane.
\item No animal that likes samurai movies is larger than Shane.
\item No dog is larger than Daisy.
\item Any animal that dislikes samurai movies is larger than Rave.
\item There is an animal that is between Rave and Shane in size.
\item There is no dog that is between Rave and Shane in size.
\item No dog is larger than itself.
\item Every dog is larger than some dog.
\item There is an animal that is smaller than every dog.
\item If there is an animal that is larger than any dog, then that animal does not like samurai movies.
\end{compactlist}

\problempart
\label{pr.QLcandies}
Using the symbolization key given, symbolize each English-language sentence into FOL.
\begin{ekey}
\item[\text{domain}] candies
\item[\atom{C}{x}] \gap{x} has chocolate in it
\item[\atom{M}{x}] \gap{x} has marzipan in it
\item[\atom{S}{x}] \gap{x} has sugar in it
\item[\atom{T}{x}] Boris has tried \gap{x}
\item[\atom{B}{x,y}] \gap{x} is better than \gap{y}
\end{ekey}
\begin{compactlist}
\item Boris has never tried any candy.
\item Marzipan is always made with sugar.
\item Some candy is sugar-free.
\item The very best candy is chocolate.
\item No candy is better than itself.
\item Boris has never tried sugar-free chocolate.
\item Boris has tried marzipan and chocolate, but never together.
%\item Boris has tried nothing that is better than sugar-free marzipan.
\item Any candy with chocolate is better than any candy without it.
\item Any candy with chocolate and marzipan is better than any candy that lacks both.
\end{compactlist}

\problempart
Using the following symbolization key:
\begin{ekey}
\item[\text{domain}] people and dishes at a potluck
\item[\atom{R}{x}] \gap{x} has run out
\item[\atom{T}{x}] \gap{x} is on the table
\item[\atom{F}{x}] \gap{x} is food
\item[\atom{P}{x}] \gap{x} is a person
\item[\atom{L}{x,y}] \gap{x} likes \gap{y}
\item[e] Eli
\item[f] Francesca
\item[g] the guacamole
\end{ekey}
symbolize the following English sentences in FOL:
\begin{compactlist}
\item All the food is on the table.
\item If the guacamole has not run out, then it is on the table.
\item Everyone likes the guacamole.
\item If anyone likes the guacamole, then Eli does.
\item Francesca only likes the dishes that have run out.
\item Francesca likes no one, and no one likes Francesca.
\item Eli likes anyone who likes the guacamole.
\item Eli likes anyone who likes the people that he likes.
\item If there is a person on the table already, then all of the food must have run out.
\end{compactlist}


\solutions
\problempart
\label{pr.FOLballet}
Using the following symbolization key:
\begin{ekey}
\item[\text{domain}] people
\item[\atom{D}{x}] \gap{x} dances ballet
\item[\atom{F}{x}] \gap{x} is female
\item[\atom{M}{x}] \gap{x} is male
\item[\atom{C}{x,y}] \gap{x} is a child of \gap{y}
\item[\atom{S}{x,y}] \gap{x} is a sibling of \gap{y}
\item[e] Elmer
\item[j] Jane
\item[p] Patrick
\end{ekey}
symbolize the following sentences in FOL:
\begin{compactlist}
\item All of Patrick's children are ballet dancers.
\item Jane is Patrick's daughter.
\item Patrick has a daughter.
\item Jane is an only child.
\item All of Patrick's sons dance ballet.
\item Patrick has no sons.
\item Jane is Elmer's niece.
\item Patrick is Elmer's brother.
\item Patrick's brothers have no children.
\item Jane is an aunt.
\item Everyone who dances ballet has a brother who also dances ballet.
\item Every woman who dances ballet is the child of someone who dances ballet.
\end{compactlist}


\chapter{Identity}
\label{sec.identity}

Consider this sentence:
\begin{enumerate}
\item\label{else1} Pavel owes money to everyone.
\end{enumerate}
Let the domain be people; this will allow us to symbolize `everyone' with a universal quantifier. Offering the symbolization key:
	\begin{ekey}
		\item[\atom{O}{x,y}] \gap{x} owes money to \gap{y}
		\item[p] Pavel
	\end{ekey}
we can symbolize \cref*{else1} by `$\forall x\, \atom{O}{p,x}$'. But this has a (perhaps) odd consequence. It requires that Pavel owes money to every member of the domain (whatever the domain may be). The domain certainly includes Pavel. So this entails that Pavel owes money to himself. And maybe we did not want to say that. Maybe we meant to leave it open if Pavel owes money to himself, something we could have expressed more precisely by using either one of the following:
	\begin{enumerate}
		\item\label{else1b} Pavel owes money to everyone \emph{else}.
		\item\label{else1c} Pavel owes money to everyone \emph{other than} Pavel.
	\end{enumerate}
But we do not have any way for dealing with the italicized words yet. The solution is to add another symbol to FOL.

\section{Adding identity}

The symbol `$=$' will be a two-place predicate. Since it will have a special meaning, we shall write it a bit differently: we put it \emph{between} two terms, rather than out front. (This should also be familiar; consider a mathematical equation like $\frac{1}{2} = 0.5$.) And the special meaning for `$=$' is given by the fact that we \emph{always} adopt the following symbolization key: 
	\begin{ekey}
		\item[{x=y}] \gap{x} is identical to \gap{y}
	\end{ekey}
This does not mean \emph{merely} that the terms to the left and right
of `$=$' refer to objects that are indistinguishable, or that all of
the same things are true of them. Rather, it means that the terms to
the left and right of `$=$' refer to \emph{one and the same object}.

To put this to use, suppose we want to symbolize this sentence:
\begin{enumerate}
\item\label{else2} Pavel is Mister Chekov.
\end{enumerate}
Let us add to our symbolization key:
	\begin{ekey}
		\item[c] Mister Chekov
	\end{ekey}
Now \cref*{else2} can be symbolized as `$p=c$'. This tells us that the names `$p$' and `$c$' both name the same thing.

We can also now deal with \cref{else1b,else1c}. Both of these sentences can be  paraphrased as `Everyone who is not Pavel is owed money by Pavel'. Paraphrasing some more, we get: `For all~$x$, if $x$ is not Pavel, then $x$ is owed money by Pavel'. Now that we are armed with our new identity symbol, we can symbolize this as `$\forall x (\enot\, {x = p} \eif \atom{O}{p,x})$'.

This last sentence contains the formula `$\enot\, {x = p}$'. That might look a bit strange, because the symbol that comes immediately after the `$\enot$' is a variable, rather than a predicate, but this is not a problem. We are simply negating the entire formula~`$x = p$'.

\section{`Only' and `except'}

In addition to sentences that use the word `else', and `other than', identity is helpful when symbolizing some sentences that contain the words `only', and~`except'. Consider:
\begin{enumerate}
\item\label{only1} Only Pavel owes money to Hikaru.
\end{enumerate}
Let `$h$' name Hikaru. Plausibly, \cref*{only1} is true if, and only if, both of the following conditions hold:
\begin{enumerate}
	\item\label{only2} Pavel owes money to Hikaru.
	\item\label{else3} No one who is not Pavel owes money to Hikaru.
\end{enumerate}
\Cref*{else3} can be symbolized by any one of:
\begin{align*}
	& \enot\exists x(\enot\, {x = p} \eand \atom{O}{x,h}),\\
	& \forall x(\enot\, {x = p} \eif \enot\atom{O}{x,h}),\\
	& \forall x(\atom{O}{x,h} \eif x = p).
\end{align*}
Thus, we can symbolize \cref*{only1} as the conjunction of one of the above with the symbolization of~\cref*{only2}, `$\atom{O}{p,h}$', or more compactly using `$\eiff$' as `$\forall x(\atom{O}{x,h} \eiff x = p)$'.

\begin{enumerate}
	\item\label{except} Everyone except Pavel owes money to Hikaru.
	\end{enumerate}
\Cref*{except} can be treated similarly, although now of course Pavel does \emph{not} owe Hikaru money. We can paraphrase it as `Everyone who is not Pavel owes Hikaru money, and Pavel does not'. Consequently, it can be symbolized as `$\forall x(\enot\, {x = p} \eif \atom{O}{x,h}) \eand \enot \atom{O}{p,h}$', or more concisely, `$\forall x(\enot\, {x = p} \eiff \atom{O}{x,h})$'. Other locutions akin to `except' such as `but' or `besides' (as used in `no one but Pavel' or `someone besides Hikaru') can be treated in similar ways.

The above treatment of so-called ``exceptives'' is not uncontentious. Some linguists think that \cref*{except} does not entail that Pavel doesn't owe Hikaru money, and so the symbolization should just be `$\forall x(\enot\, {x = p} \eif \atom{O}{x,h})$'.  There are also uses of `except' that clearly do not have that entailment, especially in mathematical writing.  For instance, you may read in a calculus textbook that ``the function~$f$ is defined everywhere except possibly at~$a$''.  That means only that for every point~$x$ other than~$a$, $f$~is defined at~$x$. It is not required that $f$ is undefined at~$a$; it's left open whether $f$ is or is not defined at~$a$.

\section{There are at least\ldots}
We can also use identity to say how many things there are of a particular kind. For example, consider these sentences:
\begin{enumerate}
\item\label{atleast1} There is at least one apple.
\item\label{atleast2} There are at least two apples.
\item\label{atleast3} There are at least three apples.
\end{enumerate}
We will use the symbolization key:
	\begin{ekey}
		\item[\atom{A}{x}] \gap{x} is an apple
	\end{ekey}
\Cref*{atleast1} does not require identity. It can be adequately symbolized by `$\exists x\, \atom{A}{x}$': There is an apple; perhaps many, but at least one.

It might be tempting to also symbolize \cref*{atleast2} without
identity, namely as `$\exists x \exists y(\atom{A}{x} \eand
\atom{A}{y})$'. Roughly, this says that there is some apple $x$ in the
domain and some apple $y$ in the domain. Since nothing precludes these
from being one and the same apple, this would be true even if there
were only one apple. In order to make sure that we are dealing with
\emph{different} apples, we need the identity predicate.
\Cref*{atleast2} needs to say that the two apples that exist are not
identical, so it can be symbolized by `$\exists x \exists
y((\atom{A}{x} \eand \atom{A}{y}) \eand \enot\, {x = y})$'.

\Cref*{atleast3} requires talking about three different apples. Now we need three existential quantifiers, and we need to make sure that each will pick out something different: 
\[
	\exists x \exists y\exists z[((\atom{A}{x} \eand \atom{A}{y}) \eand \atom{A}{z}) \eand ((\enot\, {x = y} \eand \enot\, {y = z}) \eand \enot\, {x = z})].
\]
Note that it is \emph{not} enough to use `$\enot\, {x = y} \eand \enot\, {y = z}$' to symbolize `$x$, $y$, and~$z$ are all different.' For that would be true if $x$ and $y$ were different, but $x = z$. In general, to say that $x_1$, \dots, $x_n$ are all different, we must have a conjunction of $\enot\, {x_i = x_j}$ for every different pair $i$ and~$j$.

\section{There are at most\ldots}
Now consider these sentences:
\begin{enumerate}
	\item\label{atmost1} There is at most one apple.
	\item\label{atmost2} There are at most two apples.
\end{enumerate}
\Cref*{atmost1} can be paraphrased as, `It is not the case that there
are at least \emph{two} apples'. This is just the negation of
\cref{atleast2}: 
\[
	\enot \exists x \exists y[(\atom{A}{x} \eand \atom{A}{y}) \eand \enot\, {x = y}]
\]
But \cref*{atmost1} can also be approached in another way. It means
that if you pick out an object and it's an apple, and then you pick
out an object and it's also an apple, you must have picked out the
same object both times. With this in mind, \cref*{atmost1} can also be
symbolized by:
\[
	\forall x\forall y\bigl[(\atom{A}{x} \eand \atom{A}{y}) \eif x=y\bigr]
\]
The two sentences will turn out to be logically equivalent.

Similarly, \cref*{atmost2} can be approached in two equivalent ways.
It can be paraphrased as, `It is not the case that there are
\emph{three} or more distinct apples', so we can offer:
\[
	\enot \exists x \exists y\exists z\bigl[((\atom{A}{x} \eand \atom{A}{y}) \eand \atom{A}{z}) \eand ((\enot\, {x = y} \eand \enot\, {x = z}) \eand \enot\, {y = z})\bigr]
\]
Alternatively we can read it as saying that if you pick out an apple,
and an apple, and an apple, then you will have picked out (at least)
one of these objects more than once. Thus:
\[
	\forall x\forall y\forall z\bigl[((\atom{A}{x} \eand \atom{A}{y}) \eand \atom{A}{z}) \eif ((x=y \eor x=z) \eor y=z)\bigr]
\]


\section{There are exactly\ldots}\label{sec:exactlyn}

We can now consider statements of exact numerical quantity, like:
\begin{enumerate}
\item\label{exactly1} There is exactly one apple.
\item\label{exactly2} There are exactly two apples.
\item\label{exactly3} There are exactly three apples.
\end{enumerate}
\Cref*{exactly1} can be paraphrased as, `There is \emph{at least} one
apple and there is \emph{at most} one apple'. This is just the
conjunction of \cref{atleast1} and \cref{atmost1}. So we can offer:
\[
	\exists x \atom{A}{x} \eand \forall x\forall y\bigl[(\atom{A}{x} \eand \atom{A}{y}) \eif x=y\bigr]
\]
But it is perhaps more straightforward to paraphrase \cref*{exactly1}
as, `There is a thing $x$ which is an apple, and everything which is
an apple is just $x$ itself'. Thought of in this way, we offer: 
\[
	\exists x\bigl[\atom{A}{x} \eand \forall y(\atom{A}{y} \eif x= y)\bigr]
\]
Similarly, \cref*{exactly2} may be paraphrased as, `There are \emph{at
least} two apples, and there are \emph{at most} two apples'. Thus we
could offer
\ifHTMLtarget
\[\exists x \exists y((\atom{A}{x} \eand \atom{A}{y}) \eand \enot\, {x = y}) \eand 
\forall x\forall y\forall z\bigl[((\atom{A}{x} \eand \atom{A}{y})
\eand \atom{A}{z}) \eif ((x=y \eor x=z) \eor y=z)\bigr]
\]
\else
\begin{multline*}
  \exists x \exists y((\atom{A}{x} \eand \atom{A}{y}) \eand \enot\, {x = y}) \eand {}\\
  \forall x\forall y\forall z\bigl[((\atom{A}{x} \eand \atom{A}{y}) \eand \atom{A}{z}) \eif ((x=y \eor x=z) \eor y=z)\bigr]
\end{multline*}
\fi
More efficiently, though, we can paraphrase it as `There are at least two different apples, and every apple is one of those two apples'. Then we offer:
$$\exists x\exists y\bigl[((\atom{A}{x} \eand \atom{A}{y}) \eand \enot\, {x = y}) \eand \forall z(\atom{A}{z} \eif ( x= z \eor y = z))\bigr]$$
Finally, consider these sentences:
\begin{enumerate}
\item\label{exactly2things} There are exactly two things.
\item\label{exactly2objects} There are exactly two objects.
\end{enumerate}
It might be tempting to add a predicate to our symbolization key, to symbolize the English predicate `\blank\ is a thing' or `\blank\ is an object', but this is unnecessary. Words like `thing' and `object' do not sort wheat from chaff: they apply trivially to everything, which is to say, they apply trivially to every thing. So we can symbolize either sentence with either of the following:
\begin{align*}
	& \exists x \exists y\,\enot\, {x = y} \eand  \enot \exists x \exists y \exists z ((\enot\, {x = y} \eand \enot\, {y = z}) \eand \enot\, {x = z}) \\
	& \exists x \exists y \bigl[\enot\, {x = y} \eand \forall z(x=z \eor y = z)\bigr]
\end{align*}

\practiceproblems

%\solutions
%\problempart
%\label{pr.FOLcandies}
%Using the following symbolization key:
%\begin{ekey}
%\item[\text{domain}] candies
%\item[\atom{C}{x}] \gap{x} has chocolate in it.
%\item[\atom{M}{x}] \gap{x} has marzipan in it.
%\item[\atom{S}{x}] \gap{x} has sugar in it.
%\item[\atom{T}{x}] Boris has tried \gap{x}.
%\item[\atom{B}{x,y}] \gap{x} is better than \gap{y}.
%\end{ekey}
%symbolize the following English sentences in FOL:
%\begin{compactlist}
%\item Boris has never tried any candy.
%\item Marzipan is always made with sugar.
%\item Some candy is sugar-free.
%\item The very best candy is chocolate.
%\item No candy is better than itself.
%\item Boris has never tried sugar-free chocolate.
%\item Boris has tried marzipan and chocolate, but never together.
%%\item Boris has tried nothing that is better than sugar-free marzipan.
%\item Any candy with chocolate is better than any candy without it.
%\item Any candy with chocolate and marzipan is better than any candy that lacks both.
%\end{compactlist}

\problempart 
Consider the sentence,
\begin{enumerate}
	\item\label{except2} Every officer except Pavel owes money to Hikaru.
\end{enumerate}
Symbolize this sentence, using `$\atom{F}{x}$' for `\gap{x} is an officer'.  Are you confident that your symbolization is true if, and only if, \cref*{except2} is true?  What happens if every officer owes money to Hikaru, Pavel does not, but Pavel isn't an officer?

\problempart 
Explain why:
	\begin{compactlist}
		\item `$\exists x \forall y(\atom{A}{y} \eiff x= y)$' is a good
		symbolization of `there is exactly one apple'.
		\item `$\exists x \exists y \bigl[\enot\, {x = y} \eand \forall
		z(\atom{A}{z} \eiff (x= z \eor y = z))\bigr]$' is a good
		symbolization of `there are exactly two apples'.
	\end{compactlist}


\chapter{Sentences of FOL}\label{s:FOLSentences}
We know how to represent English sentences in FOL. The time has finally come to define the notion of a \emph{sentence} of FOL.

\section{Expressions}
There are six kinds of symbols in FOL:

\begin{compactlist}
\item \textbf{Predicates:} $A$, $B$, $C$, \dots, $Z$, or 
with subscripts, as needed: $A_1$, $B_1$, $Z_1$, $A_2$, $A_{25}$,
$J_{375}$, \dots

The identity symbol `$=$' counts as a special predicate.
\item \textbf{Names:} $a$, $b$, $c$, \dots, $r$, or
with subscripts, as needed: $a_1$, $b_{224}$, $h_7$, $m_{32}$, \dots
\item \textbf{Variables:} $s$, $t$, $u$, $v$, $w$, $x$, $y$, $z$, or
with subscripts, as needed: $x_1$, $y_1$, $z_1$, $x_2$, \dots
\item \textbf{Connectives:}  $\enot$, $\eand$, $\eor$, $\eif$, $\eiff$
\item \textbf{Brackets:} ( , )
\item \textbf{Quantifiers:}  $\forall$, $\exists$
\end{compactlist}
We define an \define{expression of FOL} as any string of symbols of FOL. Take any of the symbols of FOL and write them down, in any order, and you have an expression.

\section{Terms and formulas}
\label{s:TermsFormulas}

In \cref{s:TFLSentences}, we went straight from the statement of the vocabulary of TFL to the definition of a sentence of TFL. In FOL, we will have to go via an intermediary stage: via the notion of a \define{formula}. The intuitive idea is that a formula is any sentence, or anything which can be turned into a sentence by adding quantifiers out front. But this intuitive idea will take some time to unpack.

We start by defining the notion of a term.
	\factoidbox{
		A \define{term} is any name or any variable. }
So, here are some terms:
	$$a, b, x, x_1, x_2, y, y_{254}, z$$
Next we need to define atomic formulas.
	\factoidbox{
		\begin{factoidnumber}
		\item Any sentence letter is an atomic formula.
		\item If $\metav{R}$ is an $n$-place predicate and $\metav{t}_1, \metav{t}_2, \ldots, \metav{t}_n$ are terms, then $\atom{\metav{R}}{\metav{t}_1, \metav{t}_2, \ldots, \metav{t}_n}$ is an atomic formula.
		\item If $\metav{t}_1$ and $\metav{t}_2$ are terms, then $\metav{t}_1 = \metav{t}_2$ is an atomic formula.
		\item Nothing else is an atomic formula.
		\end{factoidnumber}
	}
Note that we consider sentence letters also formulas of FOL, so every
sentence of TFL is also a formula of FOL. In \cref{ch.NDFOL}, we will
again also allow the contradiction~`$\ered$' as an atomic formula.

\newglossaryentry{term}{
  name = term,
  description = {Either a \gls{name} or a \gls{variable}}
}

\newglossaryentry{formula}{
  name = formula,
  description = {An expression of FOL built according to the inductive rules in \cref{s:TermsFormulas}}
}


The use of script letters here follows the conventions laid down in \cref{s:UseMention}. So, `$\metav{R}$' is not itself a predicate of FOL. Rather, it is a symbol of our metalanguage (augmented English) that we use to talk about any predicate of FOL. Similarly, `$\metav{t}_1$' is not a term of FOL, but a symbol of the metalanguage that we can use to talk about any term of FOL. So, where `$F$' is a one-place predicate, `$G$' is a three-place predicate, and `$S$' is a six-place predicate, here are some atomic formulas:
	\begin{center}
		\begin{tabular}{l@{\qquad\qquad}l}
		$D$ & $\atom{F}{a}$\\
		$x = a$ & $\atom{G}{x,a,y}$\\
		$a = b$ & $\atom{G}{a,a,a}$\\
		$\atom{F}{x}$ & $\atom{S}{x_1, x_2, a, b, y, x_1}$
	\end{tabular}
\end{center}
Once we know what atomic formulas are, we can offer recursion clauses to define arbitrary formulas. The first few clauses are exactly analogous to those in the definition of `sentence of TFL'.
	\factoidbox{
	\begin{factoidnumber}
		\item Every atomic formula is a formula.
		\item If \metav{A} is a formula, then $\enot\metav{A}$ is a formula.
		\item If \metav{A} and \metav{B} are formulas, then $(\metav{A}\eand\metav{B})$ is a formula.
		\item If \metav{A} and \metav{B} are formulas, then $(\metav{A}\eor\metav{B})$ is a formula.
		\item If \metav{A} and \metav{B} are formulas, then $(\metav{A}\eif\metav{B})$ is a formula.
		\item If \metav{A} and \metav{B} are formulas, then $(\metav{A}\eiff\metav{B})$ is a formula.
		\item\label{formForall} If \metav{A} is a formula and \metav{x} is a variable,
		% \metav{A} contains at least one occurrence of \metav{x}, and \metav{A} contains neither $\forall \metav{x}$ nor $\exists \metav{x}$, 
		then $\forall\metav{x}\,\metav{A}$ is a formula.
		\item\label{formExists} If \metav{A} is a formula and \metav{x} is a variable, 
		%\metav{A} contains at least one occurrence of \metav{x}, and \metav{A} contains neither $\forall \metav{x}$ nor $\exists \metav{x}$, 
		then $\exists\metav{x}\,\metav{A}$ is a formula.
		\item Nothing else is a formula.
	\end{factoidnumber}
	}
So, assuming again that `$F$' is a one-place predicate, `$G$' is a three-place predicate and `$S$' is a six place-predicate, here are some formulas you can build this way:
	\begin{align*}
		& \atom{F}{x}\\
		& \atom{G}{a,y,z}\\
		& \atom{S}{y,z,y,a,y,x}\\
		(\atom{G}{a,y,z} \eif \,& \atom{S}{y,z,y,a,y,x})\\
		\forall z (\atom{G}{a,y,z} \eif \,& \atom{S}{y,z,y,a,y,x})\\
		\atom{F}{x} \eand \forall z (\atom{G}{a,y,z} \eif \,& \atom{S}{y,z,y,a,y,x})\\
		\exists y (\atom{F}{x} \eand \forall z (\atom{G}{a,y,z} \eif \,& \atom{S}{y,z,y,a,y,x}))\\
		\forall x \exists y (\atom{F}{x} \eand \forall z (\atom{G}{a,y,z} \eif \,& \atom{S}{y,z,y,a,y,x}))
	\end{align*}
%However, this is \emph{not} a formula:
%	\begin{center}
%		$\forall x \exists x\, \atom{G}{x,x,x}$
%	\end{center}
%Certainly `$\atom{G}{x,x,x}$' is a formula, and `$\exists x\, \atom{G}{x,x,x}$' is therefore also a formula. But we cannot form a new formula by putting `$\forall x$' at the front. This violates the constraints on clause 7 of our inductive definition: `$\exists x\, \atom{G}{x,x,x}$' contains at least one occurrence of `$x$', but it already contains `$\exists x$'.

%These constraints have the effect of ensuring that variables only serve one master at any one time (see \cref{s:MultipleGenerality}). In fact, w
We can now give a formal definition of scope, which incorporates the definition of the scope of a quantifier. Here we follow the case of TFL, though we note that a logical operator can be either a connective or a quantifier:
	\factoidbox{
		\begin{factoiditems}
		\item The \define{main logical operator} in a formula is the
		operator that was introduced most recently, when that formula was
		constructed using the recursion rules.
		\item The \define{scope} of a logical operator in a formula is the
	subformula for which that operator is the main logical operator. 
\end{factoiditems}}
So we can graphically illustrate the scope of the quantifiers in the
preceding example thus:
$$\overbrace{\forall x \overbrace{\exists y (\atom{F}{x} \eiff
\overbrace{\forall z (\atom{G}{a,y,z} \eif
\atom{S}{y,z,y,a,y,x})}^{\text{scope of `}\forall
z\text{'}}}^{\text{scope of `}\exists y\text{'}})}^{\text{scope of
`$\forall x$'}}$$

\newglossaryentry{main logical operator}{
  name = main logical operator,
  description = {The operator used last in the construction of a \gls{sentence of TFL} or a \gls{formula} of FOL}
}

\newglossaryentry{scope}{
  name = scope,
  description = {The subformula of a \gls{sentence of TFL} or a \gls{formula} of FOL for which the \gls{main logical operator} is the operator}
}

\section{Sentences and free variables}\label{s:FreeVars}

Recall that we are largely concerned in logic with assertoric
sentences: sentences that can be either true or false. Many formulas
are not sentences. Consider the following symbolization key:
	\begin{ekey}
		\item[\text{domain}] people
		\item[\atom{L}{x,y}] \gap{x} loves \gap{y}
		\item[b] Boris
	\end{ekey}
Consider the atomic formula `$\atom{L}{z,z}$'. All atomic formulas are
formulas, so `$\atom{L}{z,z}$' is a formula, but can it be true or
false? You might think that it will be true just in case the person
named by `$z$' loves themself, in the same way that `$\atom{L}{b,b}$'
is true just in case Boris (the person named by `$b$') loves himself.
\emph{However, `$z$' is a variable, and does not name anyone or any
thing.}

Of course, if we put an existential quantifier out front, obtaining
`$\exists z\,\atom{L}{z,z}$', then this would be true \ifeff{} someone
loves themselves. Equally, if we wrote `$\forall z\, \atom{L}{z,z}$',
this would be true \ifeff{} everyone loves themselves. The point is
that we need a quantifier to tell us how to deal with a variable.

Let's make this idea precise.
	\factoidbox{
		An occurrence of a variable \metav{x} is \define{bound} \ifeff{} it falls within the scope of either $\forall\metav{x}$ or~$\exists\metav{x}$. An occurrence of a variable which is not bound is \define{free}.}

\newglossaryentry{bound variable}{
  name = bound variable,
  description = {An occurrence of a variable in a \gls{formula} which is in the scope of a quantifier followed by the same variable}
}

\newglossaryentry{free variable}{
  name = free variable,
  description = {An occurrence of a variable in a \gls{formula} which is not a \gls{bound variable}}
}

For example, consider the formula
	$$(\forall x(\atom{E}{x} \eor \atom{D}{y}) \eif \exists z(\atom{E}{x} \eif \atom{L}{z,x}))$$
The scope of the universal quantifier `$\forall x$' is `$\forall x (\atom{E}{x} \eor \atom{D}{y})$', so the first `$x$' is bound by the universal quantifier. However, the second and third occurrence of `$x$' are free. Equally, the `$y$' is free. The scope of the existential quantifier `$\exists z$' is `$\exists z(\atom{E}{x} \eif \atom{L}{z,x})$', so `$z$' is bound.

Formulas that do not have free variables, i.e., every variable is
bound, are called \define{sentences}. Only sentences can be true or false.
\factoidbox{	
		A \define{sentence} of FOL is any formula of FOL that contains no free variables.
	}

\newglossaryentry{sentence of FOL}{
	name = sentence (of FOL),
	text = sentence of FOL,
	plural = sentences of FOL,
	description = {A \gls{formula} of FOL which has no \glspl{free variable}}
}

Note that in the formation rules given in \cref{s:TermsFormulas},
specifically in the clauses for $\forall$ and~$\exists$, we did not
require that the variable~$\metav{x}$ in $\forall\metav{x}\,
\metav{A}$ and $\exists\metav{x}\,\metav{A}$ does not \emph{already}
occur as a bound variable in~$\metav{A}$. So we allow, for instance,
`$\exists x(\forall x\,\atom{F}{x} \eif \atom{G}{x})$' as a sentence. We
won't be meeting such special cases often, since we can always rename
bound variables to obtain an equivalent sentence that avoids such
weirdness; in this case, e.g., `$\exists x(\forall y\,\atom{F}{y} \eif
\atom{G}{x})$'. But since it is allowed to happen, we should clarify
which quantifiers bind which variables: An occurrence of a bound
variable is always bound by the quantifier with the same matching
variable that is closest to it. By ``closest'' we mean: it is the main
logical operator of the smallest subformula into the scope of which
the occurrence of the variable falls. So in our previous example,
`$\forall x$' binds the occurrence of~$x$ in `$\atom{F}{x}$', since it is
the main logical operator of `$\forall x\,\atom{F}{x}$'. The `$\forall x$'
does not bind the `$x$' in~`$\atom{G}{x}$', however. The `$\exists x$' does
not bind the `$x$' in `$\atom{F}{x}$', but it does bind the `$x$'
in~`$\atom{G}{x}$'.

The rules also do not require that the bound variable~$\metav{x}$
\emph{actually occurs} free in $\metav{A}$ (or even at all). So our
rules also allow `$\exists x\,\atom{F}{a}$' and `$\exists x\forall
x\,\atom{G}{x}$' as formulas. In each case, the `$\exists x$' does not
bind~`$x$', because `$x$' does not occur free in `$\atom{F}{a}$' or in
`$\forall x\,\atom{G}{x}$'. In the former case, `$x$' does not occur at
all; in the latter, it is already bound by~`$\forall x$'. This is called
\define{vacuous quantification}. It rarely appears in practice, and
isn't harmful. The two sentences are equivalent to the sentences
without the `$\exists x$', i.e., to `$\atom{F}{a}$' and~`$\forall
x\,\atom{G}{x}$'.

\section{Bracketing conventions}

We will adopt the same notational conventions governing brackets that we did for TFL (see \cref{s:TFLSentences,s:MoreBracketingConventions}.) First, we may omit the outermost brackets of a formula.  Second, we may use square brackets, `[' and `]', in place of brackets to increase the readability of formulas.

Sentences of FOL used in our examples can become quite cumbersome, and so we also introduce a convention to deal with conjunctions and disjunctions of more than two sentences. We stipulate that $A_1 \land A_2 \land \dots \land A_n$ and $A_1 \lor A_2 \lor \dots \lor A_n$ are to be interpreted as, respectively:
\begin{earg}
	\item[] $(\dots(A_1 \land A_2) \land \dots \land A_n)$
	\item[] $(\dots(A_1 \lor A_2) \lor \dots \lor A_n)$
\end{earg}
In practice, this just means that you are allowed to leave out parentheses in long conjunctions and disjunctions. But remember that (unless they are the outermost parentheses of the sentence) you must still enclose the entire conjunction or disjunction in parentheses. Also, you cannot mix conjunctions and disjunctions with each other or with other connectives. So the following are still not allowed, and would be ambiguous if they were:
\begin{earg}
	\item[] $A \lor B \land C \land D$
	\item[] $B \lor C \eif D$
\end{earg}

\section{Superscripts on predicates}
Above, we said that an $n$-place predicate followed by $n$ terms is an atomic formula. But there is a small issue with this definition: the symbols we use for predicates do not, themselves, indicate how many places the predicate has. Indeed, in some places in this book, we have used the letter `$G$' as a one-place predicate; in other places we have used it as a three-place predicate. So, unless we state explicitly whether we want to use `$G$' as a one-place predicate or as a three place predicate, it is \emph{indeterminate} whether `$\atom{G}{a}$' is an atomic formula or not.

There is an easy way to avoid this, which many books adopt. Instead of saying that our predicates are just capital letters (with numerical subscripts as necessary), we could say that they are capital letters \emph{with numerical superscripts} (and with numerical subscripts as necessary). The purpose of the superscript would be to say explicitly how many places the predicate has. On this approach, `$G^1$' would be a one-place predicate, and `$G^3$' would be an (entirely different) three places predicate. They would need to have different entries in any symbolisation key. And `$\atom{G^1}{a}$' would be an atomic formula, whereas `$\atom{G^3}{a}$' would not; likewise `$\atom{G^3}{a,b,c}$' would be an atomic formula, and `$\atom{G^1}{a,b,c}$' would not.

So, we \emph{could} add superscripts to all our predicates. This would have the advantage of making certain things completely explicit. However, it would have the disadvantage of making our formulas much harder to read; the superscripts would distract the eye. So, we will not bother to make this change. Our predicates will remain \emph{without} superscripts.  (And, in practice, any book which includes superscripts almost immediately stops including them!)

However, this leaves open a possibility of ambiguity. So, when any ambiguity could arise---in practice, very rarely---you should say, explicitly, how many places your predicate(s) have.

\practiceproblems
\problempart
\label{pr.freeFOL}
Identify which variables are bound and which are free.
\begin{compactlist}
\item $\exists x\, \atom{L}{x,y} \eand \forall y\, \atom{L}{y,x}$
\item $\forall x\, \atom{A}{x} \eand \atom{B}{x}$
\item $\forall x (\atom{A}{x} \eand \atom{B}{x}) \eand \forall y(\atom{C}{x} \eand \atom{D}{y})$
\item $\forall x\exists y[\atom{R}{x,y} \eif (\atom{J}{z} \eand \atom{K}{x})] \eor \atom{R}{y,x}$
\item $\forall x_1(\atom{M}{x_2} \eiff \atom{L}{x_2,x_1}) \eand \exists x_2\,\atom{L}{x_3,x_2}$
\end{compactlist}


\chapter{Definite descriptions}\label{subsec.defdesc}
Consider sentences like:
	\begin{enumerate}
		\item\label{traitor1} Nick is the traitor.
		\item\label{traitor2} The traitor went to Cambridge.
		\item\label{traitor3} The traitor is the deputy.
	\end{enumerate}
These are definite descriptions: they are meant to pick out a \emph{unique} object. They should be contrasted with \emph{indefinite} descriptions, such as `Nick  is \emph{a} traitor'. They should equally be contrasted with \emph{generics}, such as `\emph{The} whale is a mammal' (when it's inappropriate to ask \emph{which} whale). The question we face is: How should we deal with definite descriptions in FOL?


\section{Treating definite descriptions as terms}
One option would be to introduce new names whenever we come across a definite description. This is probably not a great idea. We know that \emph{the} traitor---whoever it is---is indeed \emph{a} traitor. We want to preserve that information in our symbolization.

A second option would be to use a \emph{new} definite description operator, such as `$\maththe$'. The idea would be to symbolize `the $F$' as `$\maththe x\,\atom{F}{x}$'  (think `the $x$ such that $\atom{F}{x}$'); or to symbolize `the $G$' as `$\maththe x\,\atom{G}{x}$', etc. Expressions of the form $\maththe \metav{x}\, \atom{\metav{A}}{\metav{x}}$ would then behave like names. If we were to follow this path, we could use the following symbolization key:
	\begin{ekey}
		\item[\text{domain}] people
		\item[\atom{T}{x}] \gap{x} is a traitor
		\item[\atom{D}{x}] \gap{x} is a deputy
		\item[\atom{C}{x}] \gap{x} went to Cambridge
		\item[n] Nick
	\end{ekey}
Then, we could symbolize \cref{traitor1} with `$n = \maththe x\, \atom{T}{x}$', \cref{traitor2} with `$\atom{C}{\maththe x\,\atom{T}{x}}$', and \cref{traitor3} with `$\maththe x\, \atom{T}{x} = \maththe x\, \atom{D}{x}$'.

However, it would be nice if we didn't have to add a new symbol to FOL. And we might be able to make do without one.

\section{Russell's analysis}
Bertrand Russell offered an analysis of definite descriptions. Very
briefly put, he observed that, when we say `the $F$' in the context of
a definite description, our aim is to pick out the \emph{one and only}
thing that is $F$ (in the appropriate context). Thus Russell analyzed
the notion of a definite description as follows:\footnote{Bertrand
Russell, ``On denoting'',  \emph{Mind}~14 (1905), pp.\ 479--93; also
Russell, \textit{Introduction to Mathematical Philosophy}, 1919,
London: Allen and Unwin, ch.~16.}
	\factoidbox{
		The $F$ is $G$ \textbf{\ifeff} 
		\begin{itemize}
			\item there is at least one $F$, \emph{and}
			\item there is at most one $F$, \emph{and}
			\item every $F$ is $G$.
\end{itemize}}
Note a very important feature of this analysis: \emph{`the' does not appear on the right-side of the equivalence.} Russell is aiming to provide an understanding of definite descriptions in terms that do not presuppose them.

Now, one might worry that we can say `the table is brown' without implying that there is one and only one table in the universe. But this is not (yet) a fantastic counterexample to Russell's analysis. The domain of discourse is likely to be restricted by context (e.g., to salient objects in my vicinity).

If we accept Russell's analysis of definite descriptions, then we can symbolize sentences of the form `the $F$ is $G$' using our strategy for numerical quantification in FOL. After all, we can deal with the three conjuncts on the right-hand side of Russell's analysis as follows:
	$$\exists x \atom{F}{x} \eand \forall x \forall y ((\atom{F}{x} \eand \atom{F}{y}) \eif x = y) \eand \forall x (\atom{F}{x} \eif \atom{G}{x})$$
In fact, we could express the same point rather more crisply, by recognizing that the first two conjuncts just amount to the claim that there is \emph{exactly} one $F$, and that the last conjunct tells us that that object is~$G$. So, equivalently, we could offer:
	$$\exists x \bigl[(\atom{F}{x} \eand \forall y (\atom{F}{y} \eif x = y)) \eand \atom{G}{x}\bigr]$$
Using these sorts of techniques, we can now symbolize \crefrange{traitor1}{traitor3} without using any new-fangled fancy operator, such as `$\maththe$'.

\Cref*{traitor1} is exactly like the examples we have just considered. So we would symbolize it by 
\begin{align*}
\exists x \bigl[\atom{T}{x} \eand \forall y(\atom{T}{y} \eif x = y) & \eand x = n\bigr].
\intertext{\Cref*{traitor2} poses no problems either:} 
\exists x \bigl[\atom{T}{x} \eand \forall y(\atom{T}{y} \eif x = y) & \eand \atom{C}{x}\bigr].
\end{align*}
\Cref*{traitor3} is a little trickier, because it links two definite
descriptions. But, deploying  Russell's analysis, it can be
paraphrased by `there is exactly one traitor, $x$, and there is
exactly one deputy, $y$, and $x = y$'. So we can symbolize it by:
\ifHTMLtarget
\[
	\exists x \exists y \bigl(\bigl[\atom{T}{x} \eand 
	\forall z(\atom{T}{z} \eif x = z)\bigr] \eand 
	\bigl[\atom{D}{y} \eand 
	\forall z(\atom{D}{z} \eif y = z)\bigr] \eand x = y\bigr)
\]
\else
\begin{multline*}
	\exists x \exists y \bigl(\bigl[\atom{T}{x} \eand 
	\forall z(\atom{T}{z} \eif x = z)\bigr] \eand {}\\
	\bigl[\atom{D}{y} \eand 
	\forall z(\atom{D}{z} \eif y = z)\bigr] \eand x = y\bigr)
\end{multline*}
\fi
Note that the formula `$x = y$' must fall within the scope of both quantifiers!

\section{Empty definite descriptions}

One of the nice features of Russell's analysis is that it allows us to handle \emph{empty} definite descriptions neatly.

France has no king at present. Now, if we were to introduce a name, `$k$', to name the present King of France, then everything would go wrong: remember from \cref{s:FOLBuildingBlocks} that a name must always pick out  some object in the domain, and whatever we choose as our domain, it will contain no present kings of France.

Russell's analysis neatly avoids this problem. Russell tells us to treat definite descriptions using predicates and quantifiers, instead of names. Since predicates can be empty (see \cref{s:MoreMonadic}), this means that no difficulty now arises when the definite description is empty.

Indeed, Russell's analysis helpfully highlights two ways to go wrong
in a claim involving a definite description. To adapt an example from
Stephen Neale\footnote{Stephen Neale, \textit{Descriptions},
MIT Press, 1990.} suppose Alex claims:
	\begin{enumerate}
		\item\label{kingdate} I am dating the present king of France.
	\end{enumerate}
Using the following symbolization key:
	\begin{ekey}
		\item[a] Alex
		\item[\atom{K}{x}] \gap{x} is a present king of France
		\item[\atom{D}{x,y}] \gap{x} is dating \gap{y}
	\end{ekey}
	(Note that the symbolization key speaks of \emph{a} present King of France, not \emph{the} present King of France; i.e., it employs an indefinite, rather than a definite, description.) \Cref{kingdate} would be symbolized by `$\exists x \bigl[(\atom{K}{x} \eand \forall y(\atom{K}{y} \eif  x = y)) \eand \atom{D}{a,x}\bigr]$'. Now, this can be false in (at least) two ways, corresponding to these two different sentences:
	\begin{enumerate}
		\item\label{outernegation} There is no one who is both the present King of France and such that he and Alex are dating.
		\item\label{innernegation} There is a unique present King of France, but Alex is not dating him.
	\end{enumerate}
\Cref*{outernegation} might be paraphrased by `It is not the case that: the present King of France and Alex are dating'. It will then be symbolized by `$\enot \exists x\bigl[(\atom{K}{x} \eand \forall y(\atom{K}{y} \eif  x = y)) \eand \atom{D}{a,x} \bigr]$'. We might call this \emph{outer} negation, since the negation governs the entire sentence. Note that the sentence is true if there is no present King of France.

\Cref*{innernegation} can be symbolized by `$\exists x \bigl[(\atom{K}{x} \eand \forall y(\atom{K}{y} \eif x = y)) \eand \enot \atom{D}{a,x}\bigr]$'. We might call this \emph{inner} negation, since the negation occurs within the scope of the definite description. Note that its truth requires that there is a present King of France, albeit one who is not dating Alex.
\Cref*{innernegation} entails \cref*{outernegation} (as do their
symbolizations), but not vice versa.

\section{Possessives, `both', `neither'}

We can use Russell's analysis of definite descriptions also to deal with singular possessive constructions in English.  For instance, `Smith's murderer' means something like `the person who murdered Smith', i.e., it is a disguised definite description.  On Russell's analysis, the sentence
\begin{enumerate}
	\item\label{smithsmurder} Smith's murderer is insane.
\end{enumerate}
can be false in one of three ways. It can be false because the one person who murdered Smith is not, in fact, insane. But it can also be false if the definite description is empty, namely if either no-one murdered Smith (e.g., if Smith met with an unfortunate accident) or if more than one person murdered Smith.

To symbolize sentences containing singular possessives such as `Smith's murderer' you should paraphrase them using an explicit definite description, e.g., `The person who murdered Smith is insane' and then symbolize it according to Russell's analysis. In our case, we would use the symbolization key:
\begin{ekey}
	\item[$Domain$] people
	\item[\atom{I}{x}] \gap{x} is insane 
	\item[\atom{M}{x,y}] \gap{x} murdered \gap{y}
	\item[m] Smith 
\end{ekey}
Our symbolization then reads, `$\exists x\bigl[\atom{M}{x,m} \eand \forall y(\atom{M}{y,m} \eif x=y) \eand \atom{I}{x}\bigr]$'.

Two other determiners that we can extend Russell's analysis to are
`both' and `neither'.  To say `both $F$s are~$G$' is to say that there
are exactly two $F$s, and each of them is~$G$. To say that `neither
$F$ is~$G$', is to also say that there are exactly two~$F$s, and
neither of them is~$G$. In FOL, the symbolizations would read,
respectively,
\ifHTMLtarget
\begin{align*}
	& \exists x\exists y\bigl[\atom{F}{x} \eand \atom{F}{y} \eand \enot\, {x=y} \eand 
	\forall z(\atom{F}{z} \eif (x = z \lor y = z)) \eand \atom{G}{x} \eand \atom{G}{y}\bigr]\\
	& \exists x\exists y\bigl[\atom{F}{x} \eand \atom{F}{y} \eand \enot\, {x=y} \eand 
	\forall z(\atom{F}{z} \eif (x = z \lor y = z)) \eand \enot \atom{G}{x} \eand \enot \atom{G}{y}\bigr]
\end{align*}
\else
\begin{align*}
	& \exists x\exists y\bigl[\atom{F}{x} \eand \atom{F}{y} \eand \enot\, {x=y} \eand {}\\
	& \qquad \forall z(\atom{F}{z} \eif (x = z \lor y = z)) \eand \atom{G}{x} \eand \atom{G}{y}\bigr]\\
	& \exists x\exists y\bigl[\atom{F}{x} \eand \atom{F}{y} \eand \enot\, {x=y} \eand  {}\\
	& \qquad\forall z(\atom{F}{z} \eif (x = z \lor y = z)) \eand \enot \atom{G}{x} \eand \enot \atom{G}{y}\bigr]
\end{align*}
\fi
Compare these symbolizations with the symbolizations of `exactly two $F$s are $G$s' from \cref{sec:exactlyn}, i.e., of `there are exactly two things that are both $F$ and~$G$':
\ifHTMLtarget
	\[ 
		\exists x\exists y\bigl[(\atom{F}{x} \eand \atom{G}{x}) \eand (\atom{F}{y} \eand \atom{G}{y}) \eand \enot\, {x=y} \eand 
		\forall z((\atom{F}{z} \eand \atom{G}{z}) \eif (x = z \lor y = z))\bigr]
	\]
\else
\begin{align*}
	& \exists x\exists y\bigl[(\atom{F}{x} \eand \atom{G}{x}) \eand (\atom{F}{y} \eand \atom{G}{y}) \eand \enot\, {x=y} \eand {}\\
	& \qquad\forall z((\atom{F}{z} \eand \atom{G}{z}) \eif (x = z \lor y = z))\bigr]
\end{align*}
\fi
The difference between the symbolization of this and that of `both $F$s are~$G$s' lies in the antecedent of the conditional. For `exactly two $F$s are~$G$s', we only require that there are no $F$s \emph{that are also}~$G$s other than $x$ and~$y$, whereas for `both $F$s are~$G$s', there cannot be any $F$s, whether they are~$G$s or not, other than $x$ and~$y$. In other words, `both $F$s are~$G$s' implies that exactly two~$F$s are~$G$s. However, `exactly two $F$s are~$G$s' does not imply that both $F$s are~$G$s (there might be a third~$F$ which isn't a~$G$).

\section{The adequacy of Russell's analysis}
How good is Russell's analysis of definite descriptions? This question has generated a substantial philosophical literature, but we will restrict ourselves to two observations.

One worry focusses on Russell's treatment of empty definite
descriptions. If there are no $F$s, then on Russell's analysis, both
`the $F$ is $G$' and  `the $F$ is non-$G$' are false. P.~F.\ Strawson
suggested that such sentences should not be regarded as false,
exactly, but involve \emph{presupposition failure}, and so need to be
treated as \emph{neither} true \emph{nor} false.\footnote{P.~F.\
Strawson, ``On referring'', \textit{Mind}~59 (1950), pp.\ 320--34.}

If we agree with Strawson here, we will need to revise our logic. For, in our logic, there are only two truth values (True and False), and every sentence is assigned exactly one of these truth values.

But there is room to disagree with Strawson. Strawson is appealing to some linguistic intuitions, but it is not clear that they are very robust. For example: isn't it just \emph{false}, not `gappy', that Tim is dating the present King of France?

Keith Donnellan raised a second sort of worry, which (very roughly)
can be brought out by thinking about a case of mistaken
identity.\footnote{Keith Donnellan, ``Reference and definite
descriptions'', \textit{Philosophical Review}~77 (1966), pp.\ 281--304.} Two
men stand in the corner: a very tall man drinking what looks like a
gin martini; and a very short man drinking what looks like a pint of
water. Seeing them, Malika says:
	\begin{enumerate}
		\item\label{gindrinker} The gin-drinker is very tall!
	\end{enumerate}
Russell's analysis will have us render Malika's sentence as:
	\begin{compactlist}
		\item[\ref*{gindrinker}$'$.] There is exactly one gin-drinker [in the corner], and whoever is a gin-drinker [in the corner] is very tall.
	\end{compactlist}
Now suppose that the very tall man is actually drinking \emph{water} from a martini glass; whereas the very short man is drinking a pint of (neat) gin. By Russell's analysis, Malika has said something false, but don't we want to say that Malika has said something \emph{true}? 

Again, one might wonder how clear our intuitions are on this case. We can all agree that Malika intended to pick out a particular man, and say something true of him (that he was tall). On Russell's analysis, she actually picked out a different man (the short one), and consequently said something false of him. But  maybe advocates of Russell's analysis only need to explain \emph{why} Malika's intentions were frustrated, and so why she said something false. This is easy enough to do:  Malika said something false because she had false beliefs about the men's drinks; if Malika's beliefs about the drinks had been true,  then she would have said something true.\footnote{Interested parties should read Saul Kripke, ``Speaker reference and semantic reference'', in: French et al.\ (eds.), \textit{Contemporary Perspectives in the Philosophy of Language}, University of Minnesota Press, 1977, pp.\ 6--27.}

To say much more here would lead us into deep philosophical waters. That would be no bad thing, but for now it would distract us from the immediate purpose of learning formal logic. So, for now, we will stick with Russell's analysis of definite descriptions, when it comes to putting things into FOL. It is certainly the best that we can offer, without significantly revising our logic, and it is quite defensible as an analysis.

\practiceproblems

\problempart
Using the following symbolization key:
\begin{ekey}
\item[\text{domain}] people
\item[\atom{K}{x}] \gap{x} knows the combination to the safe
\item[\atom{S}{x}] \gap{x} is a spy
\item[\atom{V}{x}] \gap{x} is a vegetarian
\item[\atom{T}{x,y}] \gap{x} trusts \gap{y}
\item[h] Hofthor
\item[i] Ingmar
\end{ekey}
symbolize the following sentences in FOL:
\begin{compactlist}
\item Hofthor trusts a vegetarian.
\item Everyone who trusts Ingmar trusts a vegetarian.
\item Everyone who trusts Ingmar trusts someone who trusts a vegetarian.
\item Only Ingmar knows the combination to the safe.
\item Ingmar trusts Hofthor, but no one else.
\item The person who knows the combination to the safe is a vegetarian.
\item The person who knows the combination to the safe is not a spy.
\end{compactlist}


\solutions
\problempart
\label{pr.FOLcards}
Using the following symbolization key:
\begin{ekey}
\item[\text{domain}] cards in a standard deck
\item[\atom{B}{x}] \gap{x} is black
\item[\atom{C}{x}] \gap{x} is a club
\item[\atom{D}{x}] \gap{x} is a deuce
\item[\atom{J}{x}] \gap{x} is a jack
\item[\atom{M}{x}] \gap{x} is a man with an axe
\item[\atom{O}{x}] \gap{x} is one-eyed
\item[\atom{W}{x}] \gap{x} is wild
\end{ekey}
symbolize each sentence in FOL:
\begin{compactlist}
\item All clubs are black cards.
\item There are no wild cards.
\item There are at least two clubs.
\item There is more than one one-eyed jack.
\item There are at most two one-eyed jacks.
\item There are two black jacks.
\item There are four deuces.
\item The deuce of clubs is a black card.
\item One-eyed jacks and the man with the axe are wild.
\item If the deuce of clubs is wild, then there is exactly one wild card.
\item The man with the axe is not a jack.
\item The deuce of clubs is not the man with the axe.
\end{compactlist}

\

\problempart Using the following symbolization key:
\begin{ekey}
\item[\text{domain}] animals in the world
\item[\atom{B}{x}] \gap{x} is in Farmer Brown's field
\item[\atom{H}{x}] \gap{x} is a horse
\item[\atom{P}{x}] \gap{x} is a Pegasus
\item[\atom{W}{x}] \gap{x} has wings
\end{ekey}
symbolize the following sentences in FOL:
\begin{compactlist}
\item There are at least three horses in the world.
\item There are at least three animals in the world.
\item There is more than one horse in Farmer Brown's field.
\item There are three horses in Farmer Brown's field.
\item There is a single winged creature in Farmer Brown's field; any other creatures in the field must be wingless.
\item The Pegasus is a winged horse.
\item The animal in Farmer Brown's field is not a horse.
\item The horse in Farmer Brown's field does not have wings.
\end{compactlist}

\problempart
In this chapter, we symbolized `Nick is the traitor' by `$\exists x (\atom{T}{x} \eand \forall y(\atom{T}{y} \eif x = y) \eand x = n)$'. Explain why these would be equally good symbolisations:
	\begin{itemize}
		\item $\atom{T}{n} \eand \forall y(\atom{T}{y} \eif n = y)$
		\item $\forall y(\atom{T}{y} \eiff y = n)$
	\end{itemize}


\chapter{Ambiguity}\label{s:ambiguityFOL}

In \cref{s:AmbiguityTFL} we discussed the fact that sentences of English can be ambiguous, and pointed out that sentences of TFL are not. One important application of this fact is that the structural ambiguity of English sentences can often, and usefully, be straightened out using different symbolizations.  One common source of ambiguity is \emph{scope ambiguity}, where the English sentence does not make it clear which logical word is supposed to be in the scope of which other. Multiple interpretations are possible.  In FOL, every connective and quantifier has a well-determined scope, and so whether or not one of them occurs in the scope of another in a given sentence of FOL is always determined.

For instance, consider the English idiom,
\begin{enumerate}
	\item\label{glitters}
	Everything that glitters is not gold.
\end{enumerate}
If we think of this sentence as of the form `every $F$ is not~$G$' where $\atom{F}{x}$ symbolizes `\gap{x} glitters' and $\atom{G}{x}$ is `\gap{x} is gold', we would symbolize it as:
\[\forall x(\atom{F}{x} \eif \enot \atom{G}{x})\]
In other words, we symbolize it the same way as we would `Nothing that glitters is gold'. But the idiom does not mean that! It means that one should not assume that just because something glitters, it is gold; not everything that appears valuable is in fact valuable.  To capture the actual meaning of the idiom, we would have to symbolize it instead as we would `Not everything that glitters is gold', i.e., in the following way:
\[\enot\forall x(\atom{F}{x} \eif \atom{G}{x})\]
Compare this with the previous symbolization: again we see that the difference in the two meanings of the ambiguous sentence lies in whether the `\enot' is in the scope of the `$\forall$' (in the first symbolization) or `$\forall$' is in the scope of `\enot' (in the second).

Of course we can alternatively symbolize the two readings using existential quantifiers as well:
\begin{align*}
	& \enot\exists x(\atom{F}{x} \eand \atom{G}{x}) \\
	& \exists x(\atom{F}{x} \eand \enot \atom{G}{x})
\end{align*}

In \cref{s:MoreMonadic} we discussed how to symbolize sentences involving `only'. Consider the sentence:
\begin{enumerate}
	\item\label{onlyamb} Only young cats are playful.
\end{enumerate} 
According to our schema, we would symbolize it this way:
\[
	\forall x(\atom{P}{x} \eif (\atom{Y}{x} \eand \atom{C}{x}))
\]
The meaning of this sentence of FOL is something like, `If an animal is playful, it is a young cat'. (Assuming that the domain is animals, of course.) This is probably not what's intended in uttering \cref*{onlyamb}, however. It's more likely that we want to say that old cats are not playful. In other words, what we mean to say is that if something is a cat and playful, it must be young. This would be symbolized as:
\[
	\forall x((\atom{C}{x} \eand \atom{P}{x}) \eif \atom{Y}{x})
\]
There is even a third reading! Suppose we're talking about young animals and their characteristics. And suppose you wanted to say that of all the young animals, only the cats are playful. You could symbolize this reading as:
\[
	\forall x((\atom{Y}{x} \eand \atom{P}{x}) \eif \atom{C}{x})
\]
Each of the last two readings can be made salient in English by placing the stress appropriately. For instance, to suggest the last reading, you would say `Only young \emph{cats} are playful', and to get the other reading you would say `Only \emph{young} cats are playful'.  The very first reading can be indicate by stressing both `young' and `cats': `Only \emph{young cats} are playful' (but not old cats, or dogs of any age).

In \cref{ss:OrderQuant,ss:SuppQuant} we discussed the importance of the order of quantifiers.  This is relevant here because, in English, the order of quantifiers is sometimes not completely determined.  When both universal (`all') and existential (`some', `a') quantifiers are  involved, this can result in scope ambiguities. Consider:
\begin{enumerate}
	\item\label{everya} Everyone went to see a movie.
\end{enumerate}
This sentence is ambiguous.  In one interpretation, it means that there is a single movie that everyone went to see. In the  other, it means that everyone went to see some movie or other, but not necessarily the same one. The two readings can be symbolized, respectively, by 
\begin{align*}
	& \exists x(\atom{M}{x} \eand \forall y(\atom{P}{y} \eif \atom{S}{y,x}))\\
	& \forall y(\atom{P}{y} \eif \exists x(\atom{M}{x} \eand \atom{S}{y,x}))
\end{align*}
We assume here that the domain contains (at least) people and movies,
and the following symbolization key:
\begin{ekey}
	\item[\atom{P}{y}] \gap{y}~is a person
	\item[\atom{M}{x}] \gap{x}~is a movie
	\item[\atom{S}{y,x}] \gap{y}~went to see~\gap{x}
\end{ekey}
In the first reading, we say that the existential quantifier has \emph{wide scope} (and its scope contains the universal quantifier, which has \emph{narrow scope}), and the other way round in the second.

In \cref{subsec.defdesc}, we encountered another scope ambiguity, arising from definite descriptions interacting with negation.  Consider Russell's own example:
\begin{enumerate}
	\item\label{thenot} The King of France is not bald.
\end{enumerate}
If the definite description has wide scope, and we are interpreting the `not' as an `inner' negation (as we said before), \cref*{thenot} is interpreted to assert the existence of a single King of France, to whom we are ascribing non-baldness. In this reading, it is symbolized as `$\exists x\bigl[\atom{K}{x} \eand \forall y(\atom{K}{y} \eif x=y)) \eand \enot \atom{B}{x}\bigr]$'. In the other reading, the `not' denies the sentence `The King of France is bald', and we would symbolize it as: `$\enot\exists x\bigl[\atom{K}{x} \eand \forall y(\atom{K}{y} \eif x=y)) \eand \atom{B}{x}\bigr]$'. In the first case, we say that the definite description has wide scope and in the second that it has narrow scope.

\practiceproblems
\problempart
Each of the following sentences is ambiguous. Provide a symbolization key for each, and symbolize all readings.
\begin{compactlist}
	\item No one likes a quitter.
	\item CSI found only red hair at the scene.
	\item Smith's murderer hasn't been arrested.
\end{compactlist}

\problempart
Russell gave the following example in his paper `On Denoting':
\begin{quote}
	I have heard of a touchy owner of a yacht to whom a guest, on first seeing it, remarked, `I thought your yacht was larger than it is'; and the owner replied, `No, my yacht is not larger than it is'.
\end{quote}
Explain what's going on.
